\chapter{搜索与相关性}
\label{chap:rs-search}

用户输入“防水通勤背包，价格不超过500元”，并不只是要求找到含有“背包”二字的商品。一款600元的防水包主题相符却超过预算；一款399元的登山包价格合格，却未必适合每天携带电脑；一篇背包选购文章可以帮助比较，但不能当作可购买的商品返回。如果系统又根据用户过去购买高价装备的记录抬高昂贵商品，个性化甚至可能抵消这次明确表达的要求。搜索需要回答的是：怎样把当前需求变成可核验的对象选择，并帮助用户完成这次任务？

前面各章已建立理解、候选、评分、展示和约束接口。显式查询把这些接口联系得更紧：理解决定采用哪些条件，检索决定哪些对象有机会参与比较，相关性排序决定哪些结果应当靠前，结果组织使用户能够识别和访问对象。发现需求不清或证据不足时，还可以询问、修正表达或补充检索。本章沿这一循环展开，但模型的列举不是一次请求必须顺次执行的步骤；每次只调用当前任务需要且预算允许的计算。

% Outline: 9.1
\section{搜索任务与相关性}
\label{sec:rs-search-task}

同样是一个搜索框，任务完成的标准可以不同。\term{已知对象查找}（known-item search）希望找到某个型号、页面或文件，首个正确对象出现得是否足够早很重要。\term{主题检索}（topical retrieval）要了解一个主题，往往需要多份互补材料。\term{交易需求}（transactional need）希望购买、预约或执行操作，结果必须具有当前可用性。\term{探索式搜索}（exploratory search）则允许用户在比较中逐渐明确问题；此时一次查询未必有唯一目标，方面覆盖与后续交互同样重要。任务类型不是永久贴在用户身上的属性，同一人可以先阅读防水材料介绍，再比较背包，随后寻找已选型号的购买页。

沿用当前目录$I_t$、候选集合$C_t$和结果列表$L_t$，另以$q_t$表示原查询。检索单位是目录中能够独立返回和评价的对象：商品可以按商品或具体规格计数，地点可以按门店计数，文档可以按整篇或片段计数。片段标识应同时关联父文档、起止位置和内容版本。若同一文章的三个重叠片段都含有答案，按片段计数可以是三次命中，按独立文档计数却只有一次；比较前必须固定单位，不能用切得更碎来制造覆盖提升。

一个候选记录至少保存对象标识$i$、来源$r$和检索分数$s^{(r)}(q_t,i)$。为支持实际访问，还可能需要标题、属性、片段位置、来源地址、更新时间和权限状态。检索分数通常只在给定方法和请求内有比较意义。它既不是商品价格，也不天然是相关概率；某个入口返回0.8、另一个返回8，不能直接断言后者更相关。

第~\ref{sec:rs-concepts}~节已区分相关性、响应和因果增量。这里进一步把相关性落实为“对象对当前需求的满足程度”，并要求标注者说明判断依据。可设等级$y(q_t,i)\in\{0,1,2,3\}$：0为无关或明确违反必要条件，1为仅同主题，2为满足主要要求但某项用途依据不足，3为充分满足本次定义的需求。这只是本章教学量表，真实任务应按评判说明另定等级。不能先选一套数字，再让不同标注者自行猜测“有用”的含义。

\Needspace{130mm}
\term{STaRK}把非结构化文本与结构化关系放入同一检索任务，例如既要商品描述合适，又要品牌关系成立。\term{FollowIR}则在保留主题的同时修改评判指令，检查原来相关的文档是否会因新条件变成不相关。它们说明相关性不能仅用主题相似替代，但并没有覆盖所有真实需求：STaRK含合成查询与补充人工查询，FollowIR依赖特定TREC集合、指令变化和人工判断范围\scite{rs-s08,rs-s16}

例如，文字高度相似的“背包防水喷剂”不是背包；可购买的499元背包若只写“防泼水”，也不能凭词面接近就判断适合用户要求的持续淋雨环境。价格上限是可验证条件，前提是明确价格口径、规格、优惠资格和时间。“适合通勤”则需要比较容量、重量、电脑夹层等用途证据，不能压成一个未经说明的布尔属性。条件数据缺失时，应保留未知并说明待核验项，而不是自动记为满足或不满足。

文本相似、点击概率与证据充分性因而需要分别判断。一个醒目标题可能吸引点击却没有提供所需信息；两段不相似的文字可能通过制造商关系共同回答问题；一段反复出现答案名称的文字也可能没有支持答案的关系。若任务要求证据，应说明每条待支持的主张需要什么材料、允许哪些来源以及材料是否仍然有效。相关性标签依赖这个任务说明，不由向量距离独自决定。

% Outline: 9.2
\section{搜索过程}
\label{sec:rs-search-process}

搜索过程的各环节以输出相接，而不是以模型名称相接。查询理解形成需求表达与条件；候选检索访问真实目录并形成可比较的集合；相关性排序形成次序；结果组织形成可识别、可操作的结果页；反馈决定更新什么或是否停止。图~\ref{fig:rs-search-cycle}中的回路表示需求和证据可以改变下一次行动，不要求每轮把已经有效的结果全部重算。

\begin{figure*}[htbp]
 \centering
 % Search outputs form a loop; dashed edges are optional updates.
\begin{tikzpicture}[
 x=1mm,y=1mm,font=\figurefont\footnotesize,
 searchstep/.style={draw=FigureModel,fill=FigureModel!6,line width=0.85pt,
   minimum width=32mm,minimum height=16mm,align=center},
 data/.style={draw=FigureInput,fill=FigureInput!6,line width=0.85pt,
   minimum width=32mm,minimum height=10mm,align=center},
 flow/.style={->,>=latex,draw=BookInk,line width=1.1pt,
   rounded corners=1.2mm},
 update/.style={flow,dashed,line width=0.85pt,draw=BookMuted}
]
 \path[use as bounding box] (0,0) rectangle (169,81);
 \node[searchstep] (query) at (20,47) {查询理解\\表达与条件};
 \node[searchstep] (retrieve) at (63,47) {候选检索\\合格对象集合};
 \node[searchstep] (rank) at (106,47) {相关性排序\\候选次序};
 \node[searchstep,draw=FigureOutput,fill=FigureOutput!6]
   (page) at (149,47) {结果组织\\可访问结果页};
 \node[data] (catalog) at (63,73) {当前对象目录};
 \node[data] (feedback) at (106,14) {用户或证据反馈};
 \draw[flow] (query.east) -- (retrieve.west);
 \draw[flow] (retrieve.east) -- (rank.west);
 \draw[flow] (rank.east) -- (page.west);
 \draw[flow] (catalog.south) -- (retrieve.north);
 \draw[flow] (page.south) -- (149,14) -- (feedback.east);
 \draw[update] (feedback.west) -- (20,14) -- (query.south);
 \draw[update] (63,14) -- (retrieve.south);
 \node[text=BookMuted,align=center] at (40,5) {修正需求或补充检索};
 \node[text=BookMuted,align=center] at (149,5) {任务未完成时反馈};
\end{tikzpicture}

 \caption{搜索的输入输出循环。实线表示当前检索的产物传递，虚线表示反馈触发的可选更新。候选来自当前目录；下游生成的答案还须逐项核验引用证据，不与候选标识生成混为一谈。}
 \label{fig:rs-search-cycle}
\end{figure*}

把一个600元商品排得很低，并不能修复“预算内候选集合”中混入该商品的问题；把不充分证据写成流畅摘要，也不能补回缺失的事实。诊断时需要沿这些输出定位错误。反过来，一个商品因新库存信息被排除，未必需要重训语言表示；用户把预算从500元改成800元，也未必需要重新解释“通勤背包”。保留中间产物和版本，有助于只更新受影响的环节。

% Outline: 9.2.1
\subsection{查询理解}
\label{sec:rs-search-query}

\term{查询理解}的产物应当能被后续检索执行，同时不丢掉原需求。纠错修复输入错误，切分确定词项边界，实体识别把名称关联到对象，条件识别提取价格、时间和范围，指代消解解释“这款”“更轻的”等省略。\term{查询改写}（query rewriting）把原需求补全或表达得更明确；\term{查询扩展}（query expansion）则增加可匹配的词或文本。改写与扩展可以同时出现，但增加“商务”“大容量”等词前，需要判断它是原意的表达，还是系统擅自增加的要求。语言分析和语义解析的完整机制归自然语言处理部分，第~\ref{part:natural-language-processing}~部分是相应入口。

贯穿查询可以保留为三类信息：对象是“背包”，用途是“通勤”，价格条件是当前适用口径下$p_i(t)\leq500$元；“防水”还需要对应可判断的防护等级或资料。检索词项可以包括“防水、通勤、背包”，结构化过滤使用价格，用途适配由相关性模型与展示属性共同说明。同一要求在哪个环节执行，取决于数据：可靠属性可先过滤，只有文本描述时需要后续核验；数据不足则仍是未知。第~\ref{sec:rs-decision-resources}~节的硬约束接口不能被“高相似度大概满足”代替。

历史$H_{u,t}$只提供请求前可用的补充依据。若查询仅为“背包”，近期的上下班用品浏览可以支持通勤解释；若查询明确给出500元上限，过去偏爱千元装备不能推翻这一条件。应保存原查询、所用历史范围、补出的条件及置信依据，使错误能够追溯到理解而不只是排序。第~\ref{sec:rs-search-hrnn}~节和第~\ref{sec:rs-search-keps}~节会把查询选择历史、实体消歧落实到评分；当通勤与登山两种解释将导致明显不同结果时，询问用户比默默选定一种更便于核验。

查询还可以包含参考图像与文字修改。例如“保持这个背包的轮廓，改成蓝色”同时指定主体、保留条件、修改条件以及“这个”的指向。对象多模态表示融合的是一个商品已有的图文资料；多模态查询表达的是这次请求，二者不是同一接口。背景颜色可能干扰主体判断，用户也未必逐项写出应保留的属性。第~\ref{sec:rs-search-tirg}~节与第~\ref{sec:rs-search-cirplant}~节将形成访问图像目录的组合向量，视觉识别及跨模态编码衔接第~\ref{part:image-processing}~部分。

缺少目标领域相关性标签时，还可以让语言模型补充检索表达。以下两种方法都不会凭生成文本创造真实商品或证据：一个主要对假设文档编码，另一个保留原查询并拼接伪文档。生成模型的基本计算参见\BookChapterRef{chap:transformer}{模型卷“Transformer”章}；这里关心生成后的文本如何进入搜索，以及新增内容是否改变了原要求。

\Needspace{100mm}
% Outline: 9.2.1.1
\subsubsection{HyDE：假设文档的查询表示}
\label{sec:rs-search-hyde}

短查询“防水通勤背包”与详细商品说明的表达方式相差很大，又没有领域标签可以直接微调检索器。\term{HyDE}（hypothetical document embeddings）先让语言模型写出一篇假设相关文档，再把它编码为查询向量。这样把短查询与文档的匹配，转成了假设文档与真实文档的表示比较。原方案调用已有指令语言模型与无监督对比编码器，不要求在目标集合上新增相关性训练\scite{rs-s30}

设生成器在指令$a$与查询$q$下产生$m$份假设文本$\widehat d_1,\ldots,\widehat d_m$，固定编码器$f$把文本映射到$\mathbb R^d$。真实对象文本$d_i$的向量$\bm v_i=f(d_i)$预先写入目录，请求向量与分数为
\begin{equation}
 \bm h_q=\frac1m\sum_{j=1}^{m}f(\widehat d_j),
 \qquad s(q,i)=\bm h_q^\top\bm v_i.
 \label{eq:rs-search-hyde}
\end{equation}
平均表示近似生成分布下的向量期望；原文也考察把$f(q)$作为额外一项加入平均。生成器和编码器的既有参数不在单次查询时更新，选择的是指令、生成数及调用预算。使用时按$\bm h_q$访问向量目录，输出真实片段或商品标识，而非直接输出$\widehat d_j$。

作为教学假设，两份生成文本的向量为$(0.8,0.4)$与$(0.6,0.6)$，均值为$(0.7,0.5)$。若两件对象向量为$(1,0)$和$(0,1)$，分数为0.7与0.5；直接编码原查询若得到$(0.3,0.7)$，次序则相反。这只展示表示变化怎样影响候选，不证明生成扩展更正确。两次生成、两次编码以及目录访问都要计入请求成本；可缓存的表达还需匹配查询与模型版本。

假设文本可能写出“采用某种涂层、售价399元”，这些都不是商品事实。即使向量压缩淡化了部分错误，也没有事实过滤的保证；多份假设还可能对应不同意图，平均后接近不了任何一种真实需求。实践中保留原查询的500元上限，对真实对象另作资格核验，并比较原查询路径与假设路径新增了什么。生成文本只提供匹配线索，不进入答案引用或商品承诺。

\Needspace{100mm}
% Outline: 9.2.1.2
\subsubsection{Query2doc：伪文档的查询扩展}
\label{sec:rs-search-query2doc}

词项检索也会遇到表达差异，却不能直接使用HyDE的向量。\term{Query2doc}让语言模型根据少样例提示产生伪文档$d'$，再与原查询组合。原文用四组查询—文档示例说明所需输出；针对词项检索，将原查询重复$n$次后拼接$d'$，以避免较长伪文档淹没原词，实验采用$n=5$。针对稠密检索，则使用$q$、分隔符和$d'$的一次拼接。这两种输入都显式保留原查询，但重复词语本身不构成约束保证\scite{rs-s31}

令$\operatorname{concat}$表示文本拼接，输入分别为
\[
 q_{\mathrm{lex}}=\operatorname{concat}(\underbrace{q,\ldots,q}_{n\text{次}},d'),
 \qquad
 q_{\mathrm{dense}}=\operatorname{concat}(q,\texttt{[SEP]},d').
\]
语言模型负责生成表达；倒排目录或向量目录负责取得真实对象。若直接在既有检索器上使用，单次扩展无需再训练它。原文的稠密实验还分别考察从BERT初始化、用正例与困难负例作对比训练，以及用交叉编码教师的候选分布作KL蒸馏并结合对比损失。后者改变了检索器训练，不能与仅增加查询文本的收益混算。

假设原词为“防水、通勤、背包”，生成文本含“防雨背包适合上下班，可放电脑”。采用的扩展应保留原价格条件，同时把“防雨、上下班”视为待验证的匹配入口；“电脑”只有在语境支持时才适合补充，不能自动变成硬条件。设某对象对三个原词的BM25基础贡献之和为0.7，对新增词的贡献之和为0.4，且本次实现线性计入查询词频，则重复五次的总贡献为$5\times0.7+0.4=3.9$。不重复时为1.1。若实现先把查询词去重，这种重复就不会产生该效果，必须核对具体输入接口。

HyDE通常平均假设文档的编码，Query2doc的稠密路径则编码包含原查询的整段拼接文本。前者改变查询向量的构造，后者还可改变倒排访问的词项及权重；不能只按“都调用语言模型”把两者视为相同算法。两者都应记录生成长度、生成与检索时延、新增候选覆盖及条件违反率。伪文档越长不一定越有用，额外词项增加访问量，也可能把未提出的品牌或用途带入结果。

% Outline: 9.2.2
\subsection{候选检索}
\label{sec:rs-search-retrieval}

需求表达并不是结果集合。候选检索必须在当前目录$I_t$中找到可以返回的对象，保存其标识、来源和必要分数，再交给后续比较。第~\ref{sec:rs-candidate-maintenance}~节讨论的版本一致性在此同样重要：昨天编码的文档今天可能失去访问权限，昨日399元的商品今天可能已涨价。模型判断的“相近”与当前可用性必须通过同一对象标识接上。

本节把取得对象与构建集合分别讲解。候选产生比较目录检索和标识生成；集合构建负责各来源共有的资格核验、去重和融合。这是职责划分，不是固定的先后顺序。有可靠价格字段时，可以先把访问范围限制在500元以内；属性只能从描述中抽取时，则需要取回后核验，并检查后置过滤是否使候选数量不足。

% Outline: 9.2.2.1
\subsubsection{候选产生}
\label{sec:rs-search-generation}

目录检索预先组织对象的词项或向量，请求到来后按匹配关系访问；标识生成学习从查询输出对象的标识序列，再由映射表取得实际内容。表~\ref{tab:rs-search-retrieval-interfaces}给出两类接口中的代表实现。所谓混合检索，通常组合这些已有入口并融合返回集合，不是另一个脱离目录和标识的对象来源。

\begin{table*}[htbp]
 \centering
 \begin{tabularx}{\linewidth}{lXl}
  \toprule
  实现 & 对象表示与匹配计算 & 取得对象的接口 \\
  \midrule
  BM25 & 词频、集合统计及长度归一 & 词项倒排目录 \\
  DPR & 查询与片段的单向量内积 & 稠密向量目录 \\
  SPLADE & 词表空间的学习权重内积 & 词项倒排目录 \\
  ColBERT、ColBERTv2 & 词元向量的后期交互 & 多向量目录 \\
  DSI、NCI & 查询到对象标识的条件概率 & 标识解码与对象映射 \\
  \bottomrule
 \end{tabularx}
 \caption{搜索候选产生的表示、计算与访问接口。相同对象可以同时拥有多种表示；比较时应固定检索单位与候选预算。}
 \label{tab:rs-search-retrieval-interfaces}
\end{table*}

这些实现处理的困难不同。BM25提供精确词项基线，DPR和SPLADE从不同表示空间缓解词汇不匹配，ColBERT保留更细的词元匹配，ColBERTv2处理多向量存储与监督成本。Contriever、E5和ReasonIR进一步改变训练信号，BGE-M3在一个编码器中提供多类表示。DSI与NCI则改变输出接口，TIRG和CIRPLANT改变请求中的图文组合方式。它们是可比较或组合的实现，并不构成一条必须依次运行的模型链。

文本之外的结构仍可直接参与访问。若要求“由总部位于某城的厂商生产”，需要商品、制造商和所在地之间的关系记录；出现该城市名称的商品描述不足以证明关系成立。访问结构化字段、关系与文本的成本应一并记录。向量索引和倒排压缩的工程细节归系统卷；本节关注怎样形成匹配分数和真实候选，困难负例及索引采样复用第~\ref{sec:rs-retrieval-training}~节与第~\ref{sec:rs-ance}~节。

\Needspace{100mm}
% Outline: 9.2.2.1.1
\paragraph{BM25：词项检索}
\label{sec:rs-search-bm25}

若要找精确型号“B17”，先检查哪些对象确实含有这个词，比先学习复杂表示更直接。倒排目录为每个词项$v$保存出现它的对象标识及词频；查询访问相应倒排表，再按对象合并贡献。仅计出现次数会过度奖励重复长文，所有词等权又无法区分常见词“商品”与罕见型号。\term{BM25}把词频饱和、集合区分度和长度归一纳入同一评分。其概率相关性框架有更早的研究基础，2009年的文献是系统整理，并非方法首次提出的年份\scite{rs-h01}

设集合含$N$份文档，$\operatorname{df}(v)$为含$v$的文档数，$f_{v,i}$为词项$v$在$d_i$中的次数，$\ell_i$和$\bar\ell$为当前分词口径下的文档长度与集合平均长度。本节采用非负逆文档频率变体，并对去重后的查询词集合$V_q$求和：
\begin{equation}
 \begin{split}
 \operatorname{IDF}(v)&=\log\left(1+\frac{N-\operatorname{df}(v)+0.5}
                                      {\operatorname{df}(v)+0.5}\right),\\
 s(q,i)&=\sum_{v\in V_q}\operatorname{IDF}(v)
 \frac{f_{v,i}(k_1+1)}
 {f_{v,i}+k_1(1-b+b\ell_i/\bar\ell)} .
 \end{split}
 \label{eq:rs-search-bm25}
\end{equation}
原始RSJ形式的对数内不加1，常见词可能得到负权重；这里明确所用变体，避免把实现差异藏在同一个名称下。$k_1>0$控制词频饱和速度，$b\in[0,1]$控制长度修正。$b=0$不修正长度；$b=1$使用完整长度比例。固定长度时，词频增大带来的贡献趋向$\operatorname{IDF}(v)(k_1+1)$，所以同一个词重复十次不等于贡献扩大十倍。

表~\ref{tab:rs-search-bm25-example}是一组假设的已分词商品说明，三篇均为背包，长度还包含其他词。查询中价格条件另作核验，此处只计算“防水、通勤、背包”。由表可得$N=3$、$\bar\ell=200$，三词的文档频率分别为2、1、3，逆文档频率约为0.4700、0.9808、0.1335。
\begin{table}[htbp]
 \centering
 \begin{tabular}{lrrrr}
  \toprule
  对象 & 长度 & 防水 & 通勤 & 背包 \\
  \midrule
  $A$ & 100 & 2 & 1 & 3 \\
  $B$ & 200 & 1 & 0 & 2 \\
  $C$ & 300 & 0 & 0 & 4 \\
  \bottomrule
 \end{tabular}
 \caption{BM25教学输入。数字为词元计数，不是实测商品数据。}
 \label{tab:rs-search-bm25-example}
\end{table}

取$k_1=1.2,b=0.75$，$A$的长度项为$1.2(0.25+0.75\times0.5)=0.75$，三项贡献之和约为2.2201。类似可得$B$为0.6536、$C$为0.2080；截取前两个得到$\{A,B\}$。$C$四次出现“背包”仍不能抵消用途词缺失和长度影响。若增加精确型号查询，只有含该型号的倒排表会提供对应贡献；若$C$只写“防雨、上下班”，本式又不能自行知道它们对应“防水、通勤”。

BM25不通过相关文档反向传播学习表示。集合更新后重算或维护词频、文档频率和平均长度，验证集按固定任务选择$k_1,b$及分析器设置；分词、停用词和同义词词典的改变也属于方法配置。它为快速、可追踪的文本匹配提供基线，但不能凭高分断言价格、库存、否定语义或防水等级合格。把商品说明反复写得像查询，还可能抬高词面分数，须在独立标签上检查。

\Needspace{100mm}
% Outline: 9.2.2.1.2
\paragraph{DPR：双编码片段检索}
\label{sec:rs-search-dpr}

用户问“哪种包能保护雨中携带的电脑”，相关片段可能只写“密封拉链与独立电脑仓”。词项不重合时，需要从已知相关查询—片段对学习表示。\term{DPR}（dense passage retriever）使用独立的查询编码器$f_\theta$和片段编码器$g_\phi$，原方案以两个BERT的末层分类词元表示作为向量，并用内积$s(q,i)=f_\theta(q)^\top g_\phi(d_i)$匹配\scite{rs-s01}

给定正片段$i^+$及负片段集合$N_q$，最小化
\begin{equation}
 \mathcal L_{\mathrm{DPR}}(q)=
 -\log\frac{\exp s(q,i^+)}
 {\exp s(q,i^+)+\sum_{j\in N_q}\exp s(q,j)} .
 \label{eq:rs-search-dpr}
\end{equation}
原训练比较随机负例、BM25困难负例与批内其他查询的正片段，并将批内负例和BM25困难负例结合。梯度更新两侧编码器，而不是只训练最后一个相似度层。正例在负例竞争中获得较大概率，并不等于模型学得了目录中所有相关对象；未标注但同样有用的片段仍可能成为误负例。

训练后先为目录片段计算$\bm v_i=g_\phi(d_i)$并建立向量访问结构；请求时只编码当前查询，按内积取出$k$个真实标识。假设$\bm h_q=(0.8,0.6)$，两份片段向量分别为$(0.9,0.1)$与$(0.2,0.9)$，分数为0.78和0.70。若前者是本次正例、后者是唯一负例，正例概率为$1/(1+\exp(-0.08))\approx0.5200$，损失约0.6539。小幅领先仍会产生推动分离的梯度，不因次序已经正确就停止更新。

这个二维例子不表示每个坐标都是“防水”或“通勤”。内积学到的是监督下的匹配关系，仍要从原文或结构化资料核验条件。DPR的开放域问答评价常问“前$k$片段中是否至少一份含有答案字符串”；这与找到全部相关商品、取得两跳关系证据都不同。对于总部问题，出现城市名字但没有“总部位于”关系的片段可能命中字符串，却无法独自支持答案。更换片段编码器后也必须重建相应目录表示，不能让新查询向量访问旧空间而不验证兼容性。

\Needspace{100mm}
% Outline: 9.2.2.1.3
\paragraph{SPLADE：稀疏词项扩展}
\label{sec:rs-search-splade}

想缓解“防水”与“防雨”的词汇差异，又希望沿用倒排目录，可以学习词表空间中的扩展权重。\term{SPLADE}把每个输入词元的上下文表示投影到完整词表，允许输入未出现的词获得非零值。本节采用2021年原始论文的求和池化版本，不把后续最大池化与蒸馏版本混入同一公式\scite{rs-s04}

设文本$x$有$n_x$个词元，$\bm h_a$是第$a$个位置的表示，$\bm e_v$是词表$V$中词元$v$的嵌入。经原掩码语言模型头的变换$T$得到logit
\begin{equation}
 \begin{split}
 z_{a,v}=T(\bm h_a)^\top\bm e_v+b_v,\qquad
 w_v(x)=\sum_{a=1}^{n_x}\log(1+\operatorname{ReLU}(z_{a,v})).
 \end{split}
 \label{eq:rs-search-splade-representation}
\end{equation}
非负变换使负logit对应零贡献，对数抑制个别位置的过大值；求和得到$\lvert V\rvert$维权重向量$\bm w(x)$。词表投影提供可扩展的坐标，相关性监督才决定哪些坐标有用。仅使用对数并不能指定非零项预算，稀疏性还由训练正则控制。

匹配分数为$s(q,i)=\bm w(q)^\top\bm w(d_i)$。训练沿用正例、BM25困难负例和批内负例的softmax对比损失，并加上查询与文档各自的稀疏惩罚：
\begin{equation}
 \begin{split}
 \mathcal L=\mathcal L_{\mathrm{rank}}+
 \lambda_q\mathcal R_q+\lambda_d\mathcal R_d,\qquad\\
 \mathcal R_{\mathrm{FLOPS}}
 =\sum_{v\in V}\left(\frac1B\sum_{b=1}^{B}w_v(x_b)\right)^2 .
 \end{split}
 \label{eq:rs-search-splade-objective}
\end{equation}
这里$B$为批大小，FLOPS正则惩罚某些词在整批中持续取得较高权重，是平均匹配操作数的平滑近似；原文也比较$\ell_1$正则。较大的$\lambda_q$可压缩请求的访问词数，但可能减少覆盖。编码器和词表投影随联合目标更新，验证时同时观察效果、非零项数及真实倒排访问量。

假设词表只展示“防水、防雨、背包、通勤”四个相关坐标，查询权重为$(1.2,0.6,1,0.4)$，两篇说明为$(0,1.5,0.8,0.2)$和$(0.9,0,0.7,0)$。内积分别为$0.9+0.8+0.08=1.78$和$1.08+0.7=1.78$。前者即使不出现“防水”，也可经共享的扩展词“防雨”被访问。查询四个非零项意味着访问四个倒排入口；各表长度不同，因此仅报告非零项数还不足以比较延迟。

使用时离线保存文档非零词项和权重，在线编码查询并累加共享词项贡献，输出对象标识而不是扩展出来的文字。两个对象本例同分，需要稳定的并列处理，不能把小数差异当成确定偏好。扩展词权重只解释“哪些坐标产生匹配”，不是属性的证明：一个写“不防水”的说明也可能激活“防水”，500与5000的词元关系也不构成数值比较。

\Needspace{100mm}
% Outline: 9.2.2.1.4
\paragraph{ColBERT：词元后期交互检索}
\label{sec:rs-search-colbert}

一条需求包含对象、用途和防护条件，压成一个向量可能丢失局部线索；对每个候选重新联合编码又很昂贵。\term{ColBERT}把交互推迟到独立编码之后：共享BERT分别处理查询与文档，使用不同的查询、文档标记，经低维投影和归一化保留多个上下文化词元向量。原查询还补充掩码位置作为可学习扩展，文档侧过滤标点向量，减少存储\scite{rs-s02}

令查询矩阵$\bm Q\in\mathbb R^{m\times d}$、文档矩阵$\bm D_i\in\mathbb R^{n_i\times d}$各行具有单位范数。后期交互采用
\begin{equation}
 s(q,i)=\sum_{a=1}^{m}\max_{1\leq b\leq n_i}
             \bm Q_{a,:}\bm D_{i,b,:}^{\top}.
 \label{eq:rs-search-maxsim}
\end{equation}
每个查询词元找到最相近的文档词元，再把这些最大值相加。MaxSim本身没有可学习参数；原训练将正负文档得分输入成对softmax交叉熵，更新BERT、投影及专用标记。因而不是预先计算一次通用词向量相似度就完成训练。

为复算这一步，假设仅展示三个查询词元与四个文档词元的余弦相似度：
\[
 \bm Q\bm D_i^\top=
 \begin{pmatrix}
 0.9&0.1&0.2&0.3\\
 0.2&0.8&0.4&0.1\\
 0.1&0.2&0.7&0.6
 \end{pmatrix}.
\]
逐行最大值为0.9、0.8、0.7，总分2.4。另一对象若三行最大值为0.7、0.7、0.6，总分2.0；只含这两个候选的正例概率约0.5987，交叉熵约0.5130。这些矩阵是教学输入，不声称真实词元各自独立地编码了一个属性。

文档编码可离线完成。全目录检索时，先为查询词元寻找近邻词元，按所属文档取得候选，再加载候选的完整词元表示计算MaxSim并返回标识。这里只访问部分近邻时存在截断风险，应与精确计算的结果比较。若直接为已有候选打分，则同一模型也可作为重排器；能否产生全目录候选取决于部署的访问路径，而不只取决于模型名。

多向量让局部匹配可保留，但比每对象单向量需要更多存储、传输和相似度计算。MaxSim也不是一对一对齐：三个查询词元可以同时匹配同一个文档位置，并不要求存在三个相互一致的事实。否定、主客体关系和数量条件仍可能被高局部得分掩盖。分析错误时应看各行匹配到哪里，同时核对完整句子，不能把最大值之和直接解释成“全部条件满足”。

\Needspace{100mm}
% Outline: 9.2.2.1.5
\paragraph{ColBERTv2：压缩与监督改进}
\label{sec:rs-search-colbertv2}

保留每个词元的向量后，目录开销可能成为主要限制。\term{ColBERTv2}保留式~\eqref{eq:rs-search-maxsim}，一方面压缩词元表示，另一方面改善困难负例与教师监督。这两个改变各有作用：压缩降低表示成本，监督改变表示质量，不能把二者共同得到的效果全归因于压缩\scite{rs-s05}

对词元向量$\bm v$，先选最近质心$\bm c_{\kappa}$，再量化残差$\bm r=\bm v-\bm c_{\kappa}$。目录保存质心编号$\kappa$与量化残差$\widetilde{\bm r}$，使用时恢复$\widetilde{\bm v}=\bm c_{\kappa}+\widetilde{\bm r}$。原方案每个残差坐标使用1或2比特；128维表示采用4字节质心编号，加16或32字节残差，即20或36字节，相比原16位向量的256字节更小。此处只比较单向量编码，质心表、倒排表和文档映射另有开销。

假设二维单位向量$\bm v=(0.8,0.6)$对应质心$(0.75,0.65)$，真实残差为$(0.05,-0.05)$，教学量化恢复为$(0.04,-0.04)$，得到$(0.79,0.61)$。对查询词元$(1,0)$，内积从0.80变为0.79。若另一候选对应最大内积为0.795、其余查询项相同，压缩就会交换这两个候选的次序。该例采用恢复向量直接内积以隔离量化误差；实现若再归一化，应按归一化后的值重算。

监督方面，先用已有ColBERT取回困难候选，让交叉编码教师为查询—片段打分，再构造多候选训练组。原方案用64个片段组成一组，以教师分数分布和学生分布之间的KL散度蒸馏，同时使用批内负例交叉熵；训练后刷新目录和负例，再重复一次。这样可减少把未标注相关片段强行当作负例的问题，但教师也会犯错，不能把其高分当成独立人工事实。

请求侧为各查询词元访问最近质心对应的倒排表，解压相关词元，按文档聚合近似MaxSim；再对截取的候选加载完整表示评分。原文区分局部可见词元产生的下界与完整候选评分，不声称近似访问总能保留真正最高分文档。验证至少需要固定监督比较压缩前后、固定压缩比较监督前后，并分别报告向量字节数、整目录大小、访问时延和候选指标。

\Needspace{135mm}
% Outline: 9.2.2.1.6
\paragraph{DSI：文档标识生成}
\label{sec:rs-search-dsi}

前面的方法把表示存进外部访问结构。\term{DSI}（differentiable search index）探索另一种方式：让模型记住文档内容与标识的关联，并根据查询直接生成标识。要使一个输出码能找到真实对象，必须同时学习“这个文档叫什么”与“这个问题应找什么文档”，而不能只有查询到任意数字的训练\scite{rs-s06}

为每份$d_i$确定标识序列$\bm z_i=(z_{i,1},\ldots,z_{i,\ell_i})$，同一序列到序列模型$p_\theta$完成两种任务：索引任务输入文档文本，检索任务输入相关查询。对任一输入$x$，教师强制下的目标是
\begin{equation}
 \mathcal L(x,\bm z_i)=-\sum_{a=1}^{\ell_i}
     \log p_\theta(z_{i,a}\mid x,z_{i,1},\ldots,z_{i,a-1}).
 \label{eq:rs-search-dsi-objective}
\end{equation}
总目标按采样配比混合索引样本与检索样本，更新共享编码器和解码器。原文比较多种索引方向，主要结果采用文档到标识的直接索引；交替混合任务也避免只做后续检索训练而忘记文档映射。

标识可以是每文档独立的原子词元、普通数字字符串，也可以通过文档语义表示的层次聚类形成共享前缀。原子标识需要随目录规模扩展输出空间；数字串可复用词元，却不自然带有语义邻近关系；语义前缀增加结构先验，但并不是可靠的人工类别。图~\ref{fig:rs-search-identifiers}用假设的两层码说明解码。若首位0、1概率为0.6、0.4，后续条件概率分别为$(0.7,0.3)$和$(0.4,0.6)$，四个码的概率为0.42、0.18、0.16、0.24。保留两个完整序列时顺序是00、11，而不是先把0分支的所有叶子都返回。

\begin{figure}[htbp]
 \centering
 % Conditional probabilities and document mapping in a toy identifier tree.
\begin{tikzpicture}[
 x=1mm,y=1mm,font=\figurefont\footnotesize,
 code/.style={draw=FigureModel,fill=FigureModel!6,minimum width=14mm,
   minimum height=7mm,inner sep=2pt,align=center},
 leaf/.style={code,draw=FigureOutput,fill=FigureOutput!6},
 edge/.style={->,>=latex,draw=BookInk,line width=0.9pt}
]
 \node[code] (root) at (51,47) {查询$q$};
 \node[code] (zero) at (25,29) {前缀0};
 \node[code] (one) at (77,29) {前缀1};
 \node[leaf] (a) at (10,9) {00\\$d_A$: 0.42};
 \node[leaf] (b) at (36,9) {01\\$d_B$: 0.18};
 \node[leaf] (c) at (66,9) {10\\$d_C$: 0.16};
 \node[leaf] (d) at (92,9) {11\\$d_D$: 0.24};
 \draw[edge] (root) -- node[above left] {0.6} (zero);
 \draw[edge] (root) -- node[above right] {0.4} (one);
 \draw[edge] (zero) -- node[left] {0.7} (a);
 \draw[edge] (zero) -- node[right] {0.3} (b);
 \draw[edge] (one) -- node[left] {0.4} (c);
 \draw[edge] (one) -- node[right] {0.6} (d);
\end{tikzpicture}

 \caption{教学标识树与真实文档映射。边标的是给定查询和前缀的条件概率，叶子概率由路径相乘得到；树形本身不证明对象合格。}
 \label{fig:rs-search-identifiers}
\end{figure}

推理时束搜索近似寻找高概率标识，恢复对应文档并执行截断、合法性与资格检查。若束宽只有1，第一步选择0后就丢失11；即使11是第二高概率结果，也不在输出中。生成未登记标识、重复码或已删除对象时均不能凭生成文本补造对象。标识00的概率0.42也不是文档事实正确率。

新增文档需要建立内容、标识和参数中的访问关系；重排聚类又可能改变已有前缀。只更新映射表而不处理模型所记内容，可能把旧码错误地指向新文档。DSI的静态集合实验展示了这一输出接口的可行性，不等于解决了任意动态目录的维护。与答案生成相比，它的输出是待访问标识，评价候选命中和合法率；答案是否正确、引用是否支持主张仍属于后续任务。

\Needspace{100mm}
% Outline: 9.2.2.1.7
\paragraph{Contriever：无监督对比表示检索}
\label{sec:rs-search-contriever}

DPR的正例依赖查询—片段标注。只有大量文档而没有这些标签时，可以先学习“同一文档的两种表达应相近”。\term{Contriever}从同一文本独立裁剪两个片段，并配合随机删除等增强构成正对，以其他文档的视图作负例。共享Transformer对查询和文档编码，对末层词元表示平均池化，再以内积匹配；无监督训练得到的表示可直接建目录，也可再接受检索标签微调\scite{rs-s25}

例如，假设一段说明为“这款背包采用密封拉链，可放电脑，适合雨天上下班”。两次裁剪得到“密封拉链，可放电脑”和“可放电脑，适合雨天上下班”。这是训练时的两个视图，不是两位用户的真实查询。其训练目标采用温度$\tau>0$控制的对比交叉熵：
\begin{equation}
 \mathcal L=-\log
 \frac{\exp(s(x,x^+)/\tau)}
 {\exp(s(x,x^+)/\tau)+\sum_{x^-\in N_x}\exp(s(x,x^-)/\tau)} .
 \label{eq:rs-search-contriever-objective}
\end{equation}
若正例和唯一负例分数为0.8、0.2，教学温度为0.2，则正例概率为$1/(1+\exp(-3))\approx0.9526$，损失约0.0486。减小温度会强调当前分数差，但并不会修复错误的正负例关系。

原文比较批内负例和跨批队列，后者采用MoCo：在线查询编码器经梯度更新，键编码器以$\theta_k\leftarrow\mu\theta_k+(1-\mu)\theta_q$缓慢跟随，历史键向量作为固定负例，不反传到过去批次。$\mu$为动量系数。它能扩大负例池而不要求同等大的即时批量，也引入键表示滞后的条件。部署时使用训练好的表示函数编码真实查询与目录，不再需要随机裁剪形成正对。

假设真实查询向量为$(0.6,0.8)$，通勤说明与户外说明向量为$(0.8,0.6)$、$(1,0)$，内积为0.96、0.60，先返回通勤说明。这个访问过程和DPR相似，区别主要在表示如何取得监督，而不在输出标识的类型。同源长文也可能包含不相关主题，裁剪出的片段未必互相回答问题；其他文档又可能重复同一事实。因此无监督正对只是学习假设，不能当成用户认可。跨域、跨语言及有监督微调应分别报告，不能把加入标签后的版本仍称为纯无监督检索。

\Needspace{100mm}
% Outline: 9.2.2.1.8
\paragraph{E5：弱监督配对的稠密检索}
\label{sec:rs-search-e5}

随机裁剪不一定像真实的信息需求，但网页、问答社区和论文中已有大量自然配对。\term{E5}从帖子—评论、问题—高票回答、标题—正文、论文标题—摘要等关系构造CCPairs，再进行弱监督对比预训练。这些关系提供可扩展信号，却混有跑题评论和不完整答案；原方案先用初步模型做一致性筛选，再训练表示，而不把所有网页邻接文本都视为相关\scite{rs-s26}

E5用共享Transformer及平均池化编码两侧，输入分别加\texttt{query:}与\texttt{passage:}前缀，使同一编码器知道当前角色。将输出归一化后，以余弦相似度除以温度进入对比目标，批内其他文本对的文档侧提供负例。原预训练采用大批量，温度设为0.01；前缀、归一化和批量组成都是训练接口的一部分，部署时不应随意省略。

取两个假设单位文档向量$(0.8,0.6)$与$(0.6,0.8)$，查询表示为$(1,0)$，相似度分别0.8、0.6。用教学温度0.1演示训练时，两者logit为8、6，正例概率约0.8808，损失0.1269。推理时同一查询各候选都除以同一正常数，不改变次序，因此目录可以直接按余弦访问。若输入改成没有前缀的文本，其表示可能改变，这不是事后缩放能补偿的差异。

原方案还提供有标签微调：对MS MARCO、NQ等使用困难负例及交叉编码教师，以教师—学生KL散度和硬标签对比损失共同优化；NLI数据则提供蕴含与矛盾关系。预训练检查点与微调检查点必须分开标记，且不能把自然语言矛盾标签等同商品搜索的预算违反标签。微调更新共享编码器后，要用相应版本重新编码目录。

固定同一背包目录比较时，DPR以明确查询—片段相关监督起步，Contriever以同源视图起步，E5以筛选后的自然文本配对起步；三者均可输出单向量候选。公平对照还需固定编码模型大小、训练数据可见范围、候选数和近邻搜索预算，不能把监督来源、架构和规模同时变化的结果解释成某个损失单独胜出。弱配对形成的“相近”仍需在领域标签上验证，不能直接充当商品偏好概率或全部条件成立的证据。

\Needspace{100mm}
% Outline: 9.2.2.1.9
\paragraph{BGE-M3：共享编码的多种表示检索}
\label{sec:rs-search-bge-m3}

同一搜索库可能同时有中文商品名、英文型号和长篇说明。只保留一种匹配计算，不容易兼顾词面细节与语义覆盖。\term{BGE-M3}对应原文的M3-Embedding，以多语言、多个表示功能和多种文本粒度为目标，在共享编码器上形成稠密、稀疏及多向量表示。共享的是编码计算，不是所有入口的存储与访问结构\scite{rs-s27}

给定词元隐藏状态，稠密分支取归一化的分类词元向量。稀疏分支通过可学习线性映射和ReLU为输入中实际出现的词元赋权，同词多次出现时取最大值，再对查询与文档共有词求权重乘积和。它不同于SPLADE对整个词表预测扩展权重。多向量分支投影并归一化所有词元，用MaxSim计算后除以查询长度，即取每个查询词元最大匹配的平均值。

设三路得分为$s_{\mathrm{dense}},s_{\mathrm{lex}},s_{\mathrm{multi}}$，融合分数写为
\begin{equation}
 s_{\mathrm{mix}}=\alpha_1s_{\mathrm{dense}}+
                  \alpha_2s_{\mathrm{lex}}+\alpha_3s_{\mathrm{multi}} .
 \label{eq:rs-search-bge-m3-fusion}
\end{equation}
先以大量多语言文本对训练稠密表示，再用标注与合成数据微调三路功能。每路以及融合分数各有正负例对比目标；融合分数在候选内经softmax得到教师分布，再以交叉熵指导各单路分布，形成自知识蒸馏。教师来自模型自己的组合预测，不是额外的事实标注。权重与各路损失系数控制早期不稳定分支的影响，编码器和新投影参数共同更新。

假设同一查询下两篇文档的三路得分为$(0.8,1.0,0.6)$与$(0.7,0.2,0.8)$，采用教学权重$(1,0.3,1)$，融合后为1.70与1.56。第一篇的词面优势抵消了第二篇的细粒度语义优势。若只用多向量，次序正好相反。因此融合权重必须在验证任务中选择，并说明分数尺度，而不是因为三个数都叫“相似度”就默认等权。

使用时可以由稠密与稀疏路径分别取得候选，合并后再算多向量分数；原文为控制成本，也用多向量重排稠密路径的候选，而非总是建立第三个全目录入口。此时多向量无法补回前两路未访问的对象。评价应注明实际路径，按语言、跨语言方向和文本长度分组；长文本最大输入长度只说明模型能接收多长，不保证每个局部证据都保留在稠密向量里。共享编码仍需承担词项表、多向量表示和额外交互成本。

\Needspace{155mm}
% Outline: 9.2.2.1.10
\paragraph{NCI：语义标识的生成检索}
\label{sec:rs-search-nci}

DSI中的语义标识提供共享前缀，但同一数字在不同分支里含义不同。\term{NCI}（neural corpus indexer）继续利用层次聚类生成语义标识，并让解码器按前缀调整输出分类权重。它还用从文档生成的查询及文档片段补足训练，把对象内容与标识的映射写入参数\scite{rs-s28}

原方案先用BERT表示文档，递归执行$k$均值聚类；小簇内给文档分配末位编号，从根到叶的路径形成$\bm z_i$。原实验分支数与小簇阈值都取30。这个树来自集合表示，而非人工保证的语义分类；相似文档倾向于共享前缀，但不同目录或不同聚类种子可能得到不同码。

前缀感知权重自适应解码器PAWA为位置$a$区分位置—词元嵌入，并由查询与已有前缀生成自适应输出矩阵。用统一列向量记法，可写其接口为
\begin{equation}
 \begin{split}
 \bm W_a^{\mathrm{ada}}&=A_\psi(q,z_{<a})\bm W_a,\\
 p(z_a\mid q,z_{<a})
 &=\operatorname{softmax}((\bm W_a^{\mathrm{ada}})^\top\bm h_a).
 \end{split}
 \label{eq:rs-search-nci-pawa}
\end{equation}
其中$A_\psi$由额外Transformer解码计算实现。普通解码器虽也在隐藏状态中利用前缀，这里还让输出分类权重随前缀变化；不能简化为“普通解码器不看历史”。同一个末位0，在前缀0和1下可以对应不同局部文档簇。

训练结合真实查询、DocT5Query生成查询及文档连续片段，以标识序列交叉熵更新编码器、解码器和适应模块。原文的一致性正则不是直接令两次输出概率相等，而是对同一查询采用独立dropout前向计算，让同一解码位置的两份表示构成对比正例，与批内其他表示区分，再将此项加到序列目标。合成查询决定对象能被哪些需求表达访问，生成质量与覆盖要单独检查。

沿图~\ref{fig:rs-search-identifiers}，假设查询改为“登山背包”，模型给首位0、1的概率变为0.3、0.7，1分支内末位0、1的概率为0.8、0.2，则10与11分别为0.56、0.14。查询对应的前缀与局部选择共同决定文档，而不是共享前缀就有同分。NCI使用标识前缀树约束束搜索，只允许合法后继，再映射到实际文档。约束减少非法码，却不能避免窄束截断；目录更新还会牵涉树结构、合成查询和模型参数。静态检索效果不能直接推出无需训练便可加入任意新文档。

\Needspace{135mm}
% Outline: 9.2.2.1.11
\paragraph{ReasonIR：推理需求的稠密检索}
\label{sec:rs-search-reasonir}

“与雨有关”不足以回答“每天步行通勤、电脑不能受潮、预算500元，应核对哪些防护证据”。有帮助的文档可能解释防泼水与持续淋雨的区别，而大量雨具广告只是同主题。\term{ReasonIR}把训练重点放在这种需要背景推理的相关关系：以公开检索数据为基础，补入变化长度的合成查询及相关文档，再从有知识含量的种子文档生成困难查询，并在另一轮生成看似相关的困难负例\scite{rs-s29}

把困难负例放到另一轮，是为了让生成器已有查询和正例可供比较，避免一次同时生成时只得到显然无关的短文本。困难查询需要在不展示种子文档时仍可理解，并且不能只复述某个罕见词。合成标注的正例可能不是整个目录里唯一有用的文档，所以不能将其他全部对象都自动定为负例；负例的逻辑错误也需检查。

ReasonIR-8B从Llama3.1-8B适应为检索编码器，将因果注意改为双向注意并形成单向量表示，以正例、批内负例和所构造困难负例学习余弦对比目标。训练时更新编码器，推理时分别编码查询与文档，通过向量目录取回对象。一般写作
\begin{equation}
 \mathcal L=-\log\frac{\exp(\gamma\,\cos(\bm h_q,\bm v_{i^+}))}
 {\sum_{j\in\{i^+\}\cup N_q}\exp(\gamma\,\cos(\bm h_q,\bm v_j))},
 \qquad \gamma>0 ,
 \label{eq:rs-search-reasonir-objective}
\end{equation}
用$\gamma$明确表示logit缩放，避免把温度及其倒数混用。相关性改善可能同时来自基础模型、注意方式和训练数据，需要消融区分，不能全记在“推理”这一名称下。

假设已核验的防护说明与仅谈雨季趋势的文章相似度为0.8、0.6，教学缩放$\gamma=5$时正例概率约0.7311，损失0.3133。如果训练把后者误认成正例，优化就会把次序往错误方向推；增加负例难度不能代替标签核验。商品价格仍用当前目录字段判断，不能因为某文档解释了防水原理，就把其中提及的昂贵产品当作合格候选。

原工作另考察上线时把查询改写得更长、包含推理背景，以及取回后的语言模型重排。这些是可选的额外计算，并非稠密检索器每次都要输出推理轨迹。比较原查询与改写查询时应记录生成模型、长度、时延和新增覆盖；BRIGHT等推理密集集合的变化也不能推成日常精确型号搜索的普遍优势。相关片段是否够用、下游答案是否正确仍须分开评价。

\Needspace{100mm}
% Outline: 9.2.2.1.12
\paragraph{TIRG：图像与文字修改的组合检索}
\label{sec:rs-search-tirg}

对于“保持这个背包的轮廓，改成蓝色”，单独检索参考图会偏向原颜色，单独检索“蓝色”则会丢失轮廓。\term{TIRG}（text image residual gating）把参考图像向量$\bm v$与修改文本向量$\bm t$组合：门控分支调节应保留的图像特征，残差分支补充修改量。图像和文本编码的通用机制衔接第~\ref{part:image-processing}~部分与\BookChapterRef{chap:recurrent-networks}{模型卷“循环神经网络”章}\scite{rs-s36}

将拼接向量记为$[\bm v;\bm t]$，其核心可写为
\begin{equation}
 \begin{split}
 \bm g&=\sigma(G([\bm v;\bm t]))\odot\bm v,\qquad
 \bm r=R([\bm v;\bm t]),\\
 \bm h_q&=a_g\bm g+a_r\bm r .
 \end{split}
 \label{eq:rs-search-tirg}
\end{equation}
$G$、$R$为可学习变换，原文含非线性与归一化，$a_g,a_r$也参与学习。门控不是孤立的一组权重，它已经与图像向量逐元素相乘；残差再改变组合表示。原文在全连接特征或卷积特征处尝试组合，这个接口不等同固定的空间注意力图。

训练实例含参考图、修改文字、目标图及负图。用归一化组合向量和候选图像向量的相似度$s_j$构成$K$类softmax，最小化$-\log(\exp s_+/\sum_{j=1}^K\exp s_j)$。$K=2$时为软三元组形式$\log(1+\exp(s_--s_+))$；原文也使用整批候选分类，而非全部实验都只有一张负图。梯度更新组合网络、图像及文本表示的可训练部分，目标图侧与目录的图像编码保持一致。

假设$\bm v=(1,0)$，门控后$\bm g=(0.8,0)$，残差$\bm r=(0,0.6)$，两系数为1，组合向量是单位向量$(0.8,0.6)$。同轮廓蓝包$A$、同轮廓原色包$B$、异轮廓蓝包$C$的单位向量分别为$(0.8,0.6)$、$(1,0)$、$(0,1)$，相似度为1、0.8、0.6，平方距离为0、0.4、0.8。若用$A$作正例、$B$作唯一负例，软三元组损失为$\log(1+\exp(-0.2))\approx0.5981$。

此例中只用参考图会把$B$排第一，只用教学文字向量$(0,1)$会把$C$排第一，组合才选中$A$。这些坐标仅用于说明保留与修改的共同作用，不能将真实嵌入维度直接读成颜色或轮廓。服务时候选图像向量可预计算，在线只需形成组合表示后近邻访问。Fashion200k、MIT-States与CSS的标签条件不同，不能用其中一个结果代替真实商品属性核验；背景、主体指向和多个合法目标都可能改变实际相关性。

\Needspace{100mm}
% Outline: 9.2.2.1.13
\paragraph{CIRPLANT：预训练视觉语言表示的组合检索}
\label{sec:rs-search-cirplant}

自然图片中的修改不只涉及预设颜色或材质，也可能描述角度、数量和场景关系。\term{CIRPLANT}把视觉语言预训练模型适应到\term{组合图像检索}，原方案以OSCAR初始化组合模块。它不把参考图与每张候选图联合输入，而是先把参考图和文字转换成一个可访问图像目录的查询向量\scite{rs-s37}

参考图经ResNet、学习投影和归一化得到全局图像向量$\bm v$，与修改文字词元$\bm t_1,\ldots,\bm t_m$一同进入视觉语言Transformer。原方案不使用可选的对象标签输入，也不以分类词元作为最终组合表示，而是取更新后的图像词元$\bm h_q$。目标图使用共享图像编码与投影得到$\bm v_i$，因此能离线写入向量目录；这个全局图像词元不同于把一组对象区域各自作为词元的表示。

训练拉近组合向量与目标图、拉远采样负图。令$\delta_+=\|\bm h_q-\bm v_{i^+}\|_2$、$\delta_-=\|\bm h_q-\bm v_{i^-}\|_2$，按距离越小越好的约定，软三元组最小化形式为
\begin{equation}
 \mathcal L_{\mathrm{triplet}}=\log(1+\exp(\delta_+-\delta_-)).
 \label{eq:rs-search-cirplant-objective}
\end{equation}
对每个正对随机采样负图并平均损失，微调图像编码、投影和组合计算。正距离的偏导为正、负距离的偏导为负，梯度下降因而产生正确的拉近与拉远方向。写距离损失时必须检查这一方向，不能把“相似度越大越好”的符号直接搬过来。

继续采用上一节的三个候选，假设更新后的图像词元为$(0.8,0.6)$，到$A,B,C$的欧氏距离分别为0、$\sqrt{0.4}\approx0.6325$、$\sqrt{0.8}\approx0.8944$。以$B$作负例时损失约0.4261。上线时从目录取最近图像，得到$A$的真实标识；组合模块没有生成一张不存在的蓝包图片。预训练有助于表达文字修改，却不保证把“这个”的主体判对。

同文发布的CIRR是数据与评价协议，不是CIRPLANT的另一种网络结构。其全目录与相似子集评价将在第~\ref{sec:rs-search-cirr}~节分别计算。相似背景可能使错轮廓图靠前，未标为目标的另一张图也可能符合业务需求。应检查主体、保留属性及修改属性，再决定是表示失败、标注不完整，还是查询本身不明确；商品价格与可售状态不在图像距离里，仍由下一节的资格核验承担。

% Outline: 9.2.2.2
\subsubsection{候选集合构建}
\label{sec:rs-search-candidate-set}

多条路径返回的对象需要成为同一可比较集合。资格核验回答“现在是否允许参与”，候选融合回答“哪些来源的对象保留下来”。二者会共同改变覆盖：先各取10个再过滤，可能只剩3个；先限定合格范围再各取10个，则可能保留其他本来排在第11名之后的合格对象。因此候选日志应记录入口原始名次、过滤理由、去重关系和最终来源，不能只保存融合后的一个分数。

\paragraph*{候选资格核验}
可靠的价格、库存、时间、地点与权限字段可以先限制访问范围。对贯穿查询，若当前价格字段与用户适用价格口径一致，可直接检查$p_i(t)\leq500$；若只拿到“最低199元”的标题，却不知道对应规格和优惠资格，就没有足够依据确定整件商品合格。缺字段、过期字段、含糊抽取结果都需要保留未知状态，必要时补查或暂不作为满足条件的结果承诺。

假设两个入口分别返回$\{A,B,C\}$与$\{B,D,E\}$，其中$B$超预算，$C$无库存，$E$价格未知。去重后有5个对象，但确定合格的只有$A,D$。把$E$混入“500元以内”列表会制造条件违反；把它记成已知不合格也会混淆数据缺失与真实违反。可以按产品协议补查$E$，并从各入口继续取数，直至达到候选预算或访问成本上限。合格集合不足时应诚实展示不足，而不是降低核验标准来凑数量。

地点范围依赖坐标和时间，新闻时效依赖发布、更新口径，文档权限依赖当前用户与资源状态。核验应覆盖词项、向量及标识生成所有入口；最后一条路径并不因模型“知道文档”而获得绕过权限的理由。具体约束执行复用第~\ref{sec:rs-decision-resources}~节，并在版本变化时复核。

\paragraph*{候选融合}
第~\ref{sec:rs-candidate-fusion}~节的来源配额、去重和融合在搜索中保持同一接口。按分数融合需要验证可比较的尺度；按名次融合减少尺度依赖，却丢失相邻候选分差。用倒数名次融合为例，教学常数$c=10$时，$A$在两路分别排1、3，$D$分别排4、1，其融合值为$1/11+1/13\approx0.1678$与$1/14+1/11\approx0.1623$。$A$略高不意味着它有0.1678的相关概率；路径中未出现的对象应按明示规则不贡献该路分数。

文档与片段也可以分层访问。混合层次检索先取得文档，再在这些文档中取片段，并在两个粒度上组合稀疏和稠密分数。它可以让文档主题为局部片段提供上下文，同时把片段访问限制在较小范围；一旦父文档被第一阶段漏掉，子片段就失去机会。原文在NQ、WebQ和TriviaQA等问答集合中比较准确度、存储及延迟，不能将其特定配置写成混合检索的普遍优势\scite{rs-s07}

融合验收应固定总候选数，报告每路独有命中、重复比例、过滤损失与额外访问成本。增加一条路径若只复制已有对象，可能不增加覆盖；若带来新的合格对象但挤走关键证据，也未必改善任务。融合结果只确定后续能比较什么，不能替代相关性排序与整页方面覆盖。

% Outline: 9.2.3
\subsection{相关性排序}
\label{sec:rs-search-ranking}

经过核验的候选仍可能满足需求的程度不同。排序以查询、当前候选和允许使用的上下文为输入，在同一查询内比较对象，输出服务本次需求的次序。用于比较的特征可以是精确匹配、语义匹配、属性证据或请求前历史，但训练标签必须说明究竟是相关等级、满意点击还是其他响应。对点击预测得准，并不自动表明对人工相关性也排得准。

第~\ref{chap:rs-ranking}~章的逐对象、成对与列表目标在这里分别对应独立标签、查询内偏好和同组竞争。RankNet建立分数差到偏好的连接，LambdaRank强调交换位置的指标影响，LambdaMART用树实现更新；之后的方法主要改变如何计算文本匹配，以及从谁那里取得监督。表~\ref{tab:rs-search-rankers}按接口比较，避免把不同模型当成一次搜索中的连续阶段。

\begin{table*}[htbp]
 \centering
 \begin{tabularx}{\linewidth}{lXX}
  \toprule
  方法 & 输入与输出 & 主要学习依据或计算条件 \\
  \midrule
  RankNet、LambdaRank、LambdaMART & 查询内特征，逐对象分数 & 成对偏好、指标权重或树拟合 \\
  DRMM、K-NRM、DUET & 词项交互，逐对象分数 & 相关性比较，匹配表示不同 \\
  monoBERT & 联合文本，相关性概率 & 二分类标签，逐候选编码 \\
  monoT5 & 联合文本，标签词元分数 & 标签生成目标，指定词元归一化 \\
  RankT5 & 联合文本，实值分数 & 直接施加成对或列表损失 \\
  RankGPT、RankZephyr & 候选列表，编号排列 & 指令调用或教师排列蒸馏 \\
  Rank1 & 查询与片段，推理后判定 & 教师轨迹蒸馏，额外生成预算 \\
  HRNN、KEPS & 查询与历史，个性化分数 & 满意点击或会话成对监督 \\
  \bottomrule
 \end{tabularx}
 \caption{相关性排序的输入与输出。每项比较都需固定查询分组、候选范围和标签；整列表输入或列表损失并不意味着同时决定结果素材与页面布局。}
 \label{tab:rs-search-rankers}
\end{table*}

个性化只在当前需求允许的范围内补充解释：含糊的“背包”可以借历史偏向通勤，明确的“登山背包”则不能被旧通勤记录改写。首次查询、历史缺失、跨领域和意图冲突应单独检查，并保留公共相关性回退。曝光选择与位置、信任偏差复用第~\ref{sec:rs-selection-bias}~节和第~\ref{sec:rs-position-bias}~节；阶段截断复用第~\ref{sec:rs-stage-connection}~节。重排器只能调整已进入候选的对象，漏检仍需返回候选环节定位。

\Needspace{100mm}
% Outline: 9.2.3.1
\subsubsection{RankNet：成对相关性学习}
\label{sec:rs-search-ranknet}

若标注者容易判断“$A$比$B$更符合通勤需求”，却不容易给出精确分数，可以直接学习两者次序。\term{RankNet}对同一查询内的特征$\bm x_i,\bm x_j$使用共享评分函数$s_i=f_\theta(\bm x_i)$，将分数差转为偏好概率：
\begin{equation}
 \begin{split}
 p_{i,j}&=\sigma(\beta(s_i-s_j)),\\
 \mathcal L_{i,j}&=-\bar p_{i,j}\log p_{i,j}
 -(1-\bar p_{i,j})\log(1-p_{i,j}).
 \end{split}
 \label{eq:rs-search-ranknet}
\end{equation}
$\beta>0$控制尺度，$\bar p_{i,j}$是监督偏好；等级较高者可取1，同等级可取0.5或在协议中不形成训练对。这里的偏好概率属于该对对象的训练模型，不是单件对象的点击率\scite{rs-h11}

取教学参数$\beta=1$，相关等级$y_A=3,y_B=0$，当前分数$s_A=0.2,s_B=0.4$。模型给$A$胜出的概率为$\sigma(-0.2)\approx0.4502$，正向偏好损失约0.7981。对两个分数的偏导分别为$p_{A,B}-1\approx-0.5498$和$+0.5498$，梯度下降将提高$A$、降低$B$。若假设分数可独立调整、步长0.1，一步后约为0.2550与0.3450，差距缩小但尚未反转。真实训练通过评分网络参数反传，不一定能独立移动每个候选。

推理时不必对所有对象两两比较：对候选各算一次$f_\theta$，按分数排序即可。对分数共同加一个常数不会改变任何成对概率，因此目标主要确定相对差异。预测同分时按稳定规则处理，不能把任意打破并列的顺序当成模型确信的偏好；监督有冲突时交叉熵寻求折中，而不是保证满足所有比较。

RankNet没有显式指定哪些错序靠近顶部。相同分数差且偏好相同的两对对象，会收到同样的局部更新，即使其中一对影响首位、另一对只影响列表末端。相关等级只提供高低时，3对0与2对1也可以得到同样的目标概率；要让等级跨度和位置进入更新，需要下一节的指标权重。候选遗漏、错误标签及选择偏差同样不会由成对形式自行消除。

\Needspace{100mm}
% Outline: 9.2.3.2
\subsubsection{LambdaRank：指标加权更新}
\label{sec:rs-search-lambdarank}

搜索往往更关心靠前位置。\term{LambdaRank}在成对分数差之外，计算交换两个候选位置会使指定指标改变多少，并据此调整更新强度。指标必须先固定：这里采用增益$G(y)=2^y-1$、折损$D(r)=1/\log_2(r+1)$及NDCG@$k$；$r>k$时令$D(r)=0$。理想排序的DCG记为$Z_q>0$，则交换位置的绝对变化为
\begin{equation}
 \Delta_{i,j}=\frac{|(G(y_i)-G(y_j))(D(r_i)-D(r_j))|}{Z_q}.
 \label{eq:rs-search-lambdarank-weight}
\end{equation}
该值在当前排序上计算，本轮更新时视为权重，而不是对离散位置直接求导\scite{rs-h11}

为免梯度符号混淆，以下$\lambda_i$专指希望增加分数的负梯度方向。对$y_i>y_j$，令
\begin{equation}
 \begin{split}
 a_{i,j}=\frac{\beta\Delta_{i,j}}{1+\exp(\beta(s_i-s_j))},
 \qquad\\
 \lambda_i\mathrel{+}=a_{i,j},\quad
 \lambda_j\mathrel{-}=a_{i,j}.
 \end{split}
 \label{eq:rs-search-lambdarank-update}
\end{equation}
把同一对象参与的各对贡献相加，再通过$\sum_i\lambda_i\nabla_\theta s_i$更新评分模型。高指标影响的错序获得较大权重，已经正确且分差很大的对则趋于较小更新；指标权重并不把逻辑函数变成硬交换操作。

假设当前顺序为$B,A,C$，等级依次为0、3、2，沿用$s_B=0.4,s_A=0.2,s_C=0$，截断$k=3$。理想顺序$A,C,B$的$Z_q=7+3/\log_2 3\approx8.8928$。交换$A,B$的变化约为0.2905，交换$A,C$的变化约为0.0589。后者本已正确，交换会变坏，但绝对变化仍衡量这对次序的重要性。取$\beta=1$，两对推动$A$的幅度分别约0.1597和0.0265。相同等级差也会因位置不同得到不同权重。

重新排序、计算权重、更新参数的过程不断重复。若两对象都在截断外，对这个NDCG@$k$定义交换影响为零；若查询无任何正增益，须约定跳过或采用其他处理，不能除以零。训练使用哪个$k$、等级如何映射增益、查询对如何采样，都影响结果。LambdaRank提供面向指标的局部更新构造，不能解释为已经求出离散NDCG的全局最优；偏好与评价口径变化时，权重也应重新定义。

\Needspace{100mm}
% Outline: 9.2.3.3
\subsubsection{LambdaMART：树模型排序}
\label{sec:rs-search-lambdamart}

手工特征包含BM25、属性相容度、来源和时效时，可以用树处理非线性组合。\term{LambdaMART}在每轮当前排序上计算Lambda，再用回归树拟合这些更新方向。树的分裂原理参见\BookSectionRef{sec:regression-trees}{模型卷“基于树与规则的回归”节}，加法提升参见\BookSectionRef{sec:ensemble-gradient-boosting}{模型卷“梯度提升机与GBDT”节}；这里新增的是查询内指标权重如何进入伪响应与叶值\scite{rs-h11}

记第$m$轮前评分为$F_{m-1}(\bm x_i)$。按查询分别排序并计算$\lambda_i$，同时将每对的
$\beta^2\Delta_{i,j}\rho_{i,j}(1-\rho_{i,j})$累加到曲率$h_i$，其中$\rho_{i,j}=1/(1+\exp(\beta(s_i-s_j)))$。在全训练集的特征—伪响应上拟合回归树分区$R_{m,\ell}$，固定本轮权重的Newton叶值为
\begin{equation}
 \begin{split}
 \gamma_{m,\ell}&=
 \frac{\sum_{\bm x_i\in R_{m,\ell}}\lambda_i}
      {\sum_{\bm x_i\in R_{m,\ell}}h_i},\\
 F_m(\bm x)&=F_{m-1}(\bm x)+
 \eta\sum_\ell\gamma_{m,\ell}\mathbf1[\bm x\in R_{m,\ell}].
 \end{split}
 \label{eq:rs-search-lambdamart}
\end{equation}
分母过小需按实现设置最小曲率或正则；不能对空叶或零曲率直接计算比值。树分裂可跨查询共享，但偏好对和指标归一化不能跨查询乱配。

继续上一节三对象，加入$C>B$一对后，$A,B,C$的Lambda分别约为0.1862、$-0.2607$、0.0745，曲率分别约为0.0865、0.1124、0.0551。假设一个已知特征把$A,C$分入一叶、$B$分入另一叶，叶值分别约1.8414和$-2.3189$。取$\eta=0.1$，分数更新为$A:0.3841$、$B:0.1681$、$C:0.1841$，次序变成$A,C,B$。这是一轮给定分裂的教学计算，并不声称任意特征都能找到该分裂。

推理只需将候选特征沿各树分支计算并累加叶值，无需在线读取相关等级。训练的树数、叶数、学习率、最小叶样本量等在独立验证集选择，查询和时间划分避免同一日志的近重复泄漏。特征过期、领域变化或等级定义改变时，树仍会输出数值，但数值未必有相同含义。较好的NDCG也不恢复漏检候选，更不能把一个价格违反项的“相关性扣分”当成硬约束执行。

\Needspace{100mm}
% Outline: 9.2.3.4
\subsubsection{RankGPT：搜索重排}
\label{sec:rs-search-rankgpt}

当候选文本需要联合比较，可将第~\ref{sec:rs-rankgpt}~节的编号生成接口直接用于搜索。输入必须同时给出原查询、相关性评判条件、候选文本和唯一编号，要求仅返回已有编号的完整排列。候选中的文本是待判断内容，其中出现的命令式句子不能覆盖排序任务指令。原方法主要从已有语言模型及提示取得重排能力，不需要在每次请求时训练参数\scite{rs-s14}

假设三个已核验预算与库存的候选为：编号1是符合防护和通勤用途的$A$，编号2是用途资料不足的$B$，编号3是明确偏重登山的$C$。教学提示可以写为“查询：防水通勤背包，价格不超过500元；依据给定资料判断用途与防护是否满足需求；按相关性输出全部编号”。若返回\texttt{[1] > [3] > [2]}，只需按编号恢复$A,C,B$，不能把自然语言解释中的另一个商品名追加成候选。

列表过长时使用滑窗。例如初始顺序$D,C,B,A$、窗口3、步长1，先处理尾部$C,B,A$得到$A,C,B$，列表变成$D,A,C,B$；再处理头部$D,A,C$得到$A,C,D$，最终为$A,C,D,B$。此例只是两次假设调用的恢复过程，未被同窗比较的对象可能次序不稳，多遍处理也增加调用成本。单列表、窗口大小、方向、步长和遍数须一并登记。

返回\texttt{[1] > [1] > [4]}时，重复1、遗漏2和3、非法4都是接口失败。可按预先确定的规则回退原顺序或重试，但应统计失败率，不能静默修复后只报告有效样本的指标。输入置换测试检查模型是否过度沿袭候选原名次，模型版本和温度决定复现条件。直接调用、教师排列蒸馏与训练其他学生排序器属于不同成本配置；任何重排都无法把池外的合格背包找回来。

\Needspace{100mm}
% Outline: 9.2.3.5
\subsubsection{DRMM：匹配分布的相关性评分}
\label{sec:rs-search-drmm}

两篇说明都谈“防水背包”，但查询指定的型号只在其中一篇出现。仅比较整段语义容易淡化这个差别。\term{DRMM}（deep relevance matching model）为每个查询词保存它与全部文档词的相似度分布，让精确匹配和近义匹配占据不同位置，再学习哪些查询词更重要\scite{rs-s17}

令$M_{a,b}$为查询词$a$与文档词$b$的词向量余弦相似度。把$[-1,1)$划成若干区间，并为精确匹配1单设一箱，得到每查询词的直方图$\bm h_a$。原文比较计数、归一化计数和对数计数三种形式。前馈匹配网络产生$r_a=F_\theta(\bm h_a)$，门控通过查询词表示或IDF计算
\begin{equation}
 \begin{split}
 g_a=\frac{\exp(\bm w_g^\top\bm x_a)}
 {\sum_c\exp(\bm w_g^\top\bm x_c)},\qquad
 s(q,i)=\sum_a g_a r_a .
 \end{split}
 \label{eq:rs-search-drmm}
\end{equation}
门控在查询词之间归一化，而非文档之间；它决定当前词项证据对分数的贡献。

以型号词“B17”为例，假设与一篇文档三个词的相似度为$(1,0.6,0.1)$。用$[-1,0),[0,0.5),[0.5,1),\{1\}$四箱，计数为$(0,1,1,1)$；另一篇把型号换为“B18”，假设相似度变为$(0.8,0.6,0.1)$，计数为$(0,1,2,0)$。前者保留精确实体命中，后者只有较强软匹配。若匹配网络分别输出型号贡献0.9、0.3，另一个查询词“背包”两文都贡献0.8，门控为$(0.7,0.3)$，总分为0.87与0.45。

训练对相关与不相关文档使用$\max(0,1-s_++s_-)$，本例若前者是正例，损失为0.58。更新匹配网络和门控参数；硬分箱对相似度不连续，原方案不通过此路径端到端更新词嵌入。部署时对每个候选形成直方图、得分并排序，文档分箱前仍需要与当前查询匹配。

直方图让变长文档得到固定维度统计，却丢失词出现的位置和相互关系。“B17替代B18”与“B18替代B17”可能有相同词项分布；“不防水”也仍包含精确词。原文旨在保留匹配强度的归纳假设，不代表关系和位置在所有搜索中不重要。型号核验、否定和数值仍应结合完整上下文，不能由门控高权重直接宣布满足条件。

\Needspace{100mm}
% Outline: 9.2.3.6
\subsubsection{K-NRM：核池化的软匹配评分}
\label{sec:rs-search-knrm}

DRMM的硬分箱难以把排序梯度传回词向量。\term{K-NRM}以可微的核函数替代区间计数，使“接近某种匹配强度”的程度成为连续值，再共同学习词嵌入和排序层。核函数基础参见\BookSectionRef{sec:classic-overview-common-kernels}{模型卷“常用核与表示选择”节}，这里的核用于池化局部词相似度，而不是直接构造支持向量机分类器\scite{rs-s18}

对词相似度矩阵$M_{a,b}$，第$k$个高斯核产生
\begin{equation}
 \begin{split}
 K_{a,k}&=\sum_b\exp\left(-\frac{(M_{a,b}-\mu_k)^2}{2\sigma_k^2}\right),\\
 \phi_k&=\sum_a\log K_{a,k},\qquad
 s(q,i)=\tanh(\bm w^\top\bm\phi+b_0).
 \end{split}
 \label{eq:rs-search-knrm}
\end{equation}
$\mu_k$选定相似度中心，$\sigma_k>0$决定响应宽度；实现需对很小的$K_{a,k}$稳定求对数。原文固定一组核中心和宽度，主要让词嵌入在排序监督下变化，不能误称所有核参数都由标签自由学出。

用一个查询词、两个文档词演示：文档$A$的相似度为$(1,0.5)$，$B$为$(0.5,0.5)$。取教学核$\mu=1,\sigma=0.5$，则$K_A=1+\exp(-0.5)\approx1.6065$，$K_B=2\exp(-0.5)\approx1.2131$，对数特征约0.4741和0.1931。若只有此核、权重1且偏置0，得分约0.4415和0.1908。若另加$\mu=0.5$的核，两个近义匹配又会为$B$提供不同证据；最终权重由训练决定。

原目标是查询内成对间隔损失$\max(0,1-s_++s_-)$，本例约0.7493。梯度经过排序层、核池化和余弦相似度返回词向量，改变的是面向检索的匹配空间，不必保留原通用语义近邻。服务时对已有候选计算这些软匹配统计并按分数排序；它仍需查询—文档交互，不能仅凭一个预存标量完成所有查询的打分。

核池化保留强度分布而非全部词序。高斯核的平滑性使优化可行，却不赋予模型对否定范围、主体关系或价格比较的逻辑保证。原实验从搜索点击日志形成标签，神经模型主要使用标题；与DRMM或其他方法比较时应固定文本字段、查询划分和标签构造，不能把不同点击模型产生的“相关等级”与完整人工判断混为一谈。

\Needspace{100mm}
% Outline: 9.2.3.7
\subsubsection{DUET：词项与语义的联合匹配}
\label{sec:rs-search-duet}

精确型号能强烈限定对象，近义用途又需要语义泛化，单一分布统计未必兼顾。\term{DUET}用两个互补分支共同形成网页相关性分数：局部分支处理词项精确匹配及其位置，分布式分支从字符片段表示学习连续语义，再将两个分数相加。这里的“双分支”不意味着两个分支都从一开始对全部词元做双向交互\scite{rs-s19}

局部分支构造二值矩阵$X_{a,b}=\mathbf1[q_a=d_{i,b}]$，经卷积和全连接变换得到$s_{\mathrm{local}}$。它知道哪里有相同词，却不能从这个矩阵独自得知具体词的身份。分布式分支把词转为字符$n$元片段计数，分别对查询与文档卷积；查询池化成一个向量，文档保留多个局部窗口向量，再做逐元素乘法和后续网络计算得到$s_{\mathrm{distributed}}$。最终$s=s_{\mathrm{local}}+s_{\mathrm{distributed}}$。

训练每组包含同一查询、相关文档与采样负文档，以总分在组内softmax得到正例概率，最小化负对数似然，两个分支共同更新。原工作使用人工相关性等级组织网页候选，而不是把模型名字当作偏好标签。卷积窗口、文本截断和字符片段词表都影响能保留的信息；这些不是最后简单调一个融合权重能替代的选择。

假设查询为“B17防水通勤背包”，$A$确含型号、用途描述较短，$B$没有型号但含丰富的近义用途描述。两个分支对$A$给分$(1.2,0.3)$、对$B$给分$(0.2,0.9)$，总分为1.5和1.1。若$A$是正例，二选一概率约0.5987、损失约0.5130。只用语义分支会反向排列，只用精确分支则可能漏掉用户未指定型号时的优质同义描述；联合训练决定怎样使用两类证据。

上线仍为每个候选输出一个相关性分数，再恢复对象次序。某个分支分高只能说明匹配网络找到信号，不能把它解释为完整事实证明。局部分支保留位置也不等于自动学会全部数值与否定关系，字符相近的型号还可能被语义分支拉近。原网页搜索任务中的效果不直接证明商品个性化排序有效，迁移时需要当前领域的标签、特征范围与约束核验。

\Needspace{100mm}
% Outline: 9.2.3.8
\subsubsection{monoBERT：联合编码的相关性重排}
\label{sec:rs-search-monobert}

独立编码便于预计算，但查询中的条件不能参与文档内部表示的形成。候选已经缩小时，可以接受逐对象的联合计算。\term{monoBERT}把查询和候选片段拼为\texttt{[CLS] 查询 [SEP] 片段 [SEP]}，以不同片段类型区分两侧，BERT在整个输入中联合编码，再用分类词元表示$\bm h_{q,i}$和分类头输出相关概率$p_{q,i}$。Transformer原理参见\BookChapterRef{chap:transformer}{模型卷“Transformer”章}\scite{rs-s20}

训练从预训练BERT初始化，使用候选内相关与不相关标签最小化二分类交叉熵
\begin{equation}
 \begin{split}
 \mathcal L=-\sum_{(q,i)}
  \bigl[y_{q,i}\log p_{q,i}+(1-y_{q,i})\log(1-p_{q,i})\bigr],\\
 p_{q,i}=\sigma(\bm w^\top\bm h_{q,i}+b).
 \end{split}
 \label{eq:rs-search-monobert}
\end{equation}
这里将二类softmax等价写成一个logit差的Sigmoid形式。编码器与分类头都随梯度更新，而不是只把固定通用嵌入交给一个分类器。原工作对BM25候选进行重排，查询至多保留64词元、总输入至多512词元；截断规则决定模型到底读到了什么。

假设$A$片段明确写“适合上下班携带电脑，密封拉链”，$B$片段写“适合短时防泼水，不适合持续淋雨”。原查询中的“防水”按当前任务说明要求持续淋雨，故教学二分类标签设为$(1,0)$。若logit差为$(1.5,-0.5)$，概率约为$(0.8176,0.3775)$，两项损失约0.2014和0.4741。推理按概率或logit差排序，二者在这个二分类接口下次序相同。

这种联合计算给模型机会联系“防水”与“不适合”，不保证它每次都正确处理否定。每个候选都要重新编码，即使文档相同、查询不同也不能直接复用其最终联合表示；候选数从20增到100时，应实测批处理吞吐、总延迟和截断比例。分类概率还受训练正负比例影响，是否能跨查询用于统一阈值需另作校准。

如果关键否定句被截断，错误属于输入可见范围，不应只通过调分类阈值掩盖；如果真正相关对象不在候选中，则联合编码也没有机会读到它。因此要同时报告固定候选的重排收益和端到端候选覆盖，不能将重排器描述成已经完成全目录检索。

\Needspace{100mm}
% Outline: 9.2.3.9
\subsubsection{monoT5：标签生成的相关性重排}
\label{sec:rs-search-monot5}

预训练序列生成模型同样可以判断相关性，但必须说明生成接口怎样变成排序分数。\term{monoT5}把输入写成\texttt{Query: [Q] Document: [D] Relevant:}，相关样本的目标是\texttt{true}，不相关样本是\texttt{false}。这里生成的是标签词元，不是DSI的对象标识，也不是给用户的自然语言答案\scite{rs-s21}

微调时使用目标序列的词元交叉熵，更新T5编码器和解码器；词元预测仍相对于完整词表归一化。打分时则取指定两个标签词元的logit $z_{\mathrm{true}},z_{\mathrm{false}}$，只在二者间重归一化：
\begin{equation}
 s(q,i)=\frac{\exp z_{\mathrm{true}}}
 {\exp z_{\mathrm{true}}+\exp z_{\mathrm{false}}}.
 \label{eq:rs-search-monot5}
\end{equation}
因此训练和评分的归一化范围需要区分。直接使用“true”的原始logit可能受其他词元共同平移影响，不能未经验证就取代上述分数。原文选择能各自作为单个词元的目标词，避免多词元标签的概率聚合歧义。

假设$A$的两项logit为$(2,0)$，$B$为$(4,3)$。尽管$B$的“true”logit更大，二标签归一化后$A$约为0.8808、$B$约为0.7311，正确比较的是正负标签之间的相对证据。若$A$训练时目标词元在完整词表中的概率为0.6，则对应词元损失为$-\log0.6\approx0.5108$，不能直接用评分阶段的0.8808替代这个训练概率。

使用时为每个候选构造相同模板并恢复分数，再按分数排列；不必自由生成长段解释。标签词选择、大小写及分词都属于模型配置，随意把“true”换成“很好”会改变既有预训练和微调关系。对长文可以切片后评分，但片段到文档的聚合规则需要公开，并保持与其他方法相同的检索单位。

monoBERT通过分类头输出概率，monoT5复用生成词表中的标签位置，二者都可服务逐对象重排。生成能力本身并不扩大候选范围，也不使分数天然校准。若模型输出相关标签却找不到支持原条件的句子，仍须检查标签定义、截断和证据，而不是用流畅的生成解释证明标签正确。

\Needspace{100mm}
% Outline: 9.2.3.10
\subsubsection{RankT5：排序损失的文本模型适应}
\label{sec:rs-search-rankt5}

标签生成关注每个查询—文档对该输出哪个词，排序关注同一查询下谁应高于谁。\term{RankT5}让T5直接输出实值分数，再按查询组织成对或列表目标。原文同时研究编码器—解码器与仅编码器两种结构，因而模型名并不唯一确定是否在线调用解码器\scite{rs-s22}

编码器—解码器版本对联合查询—文档输入进行一次解码计算，取未使用词元\texttt{<extra\_id\_10>}的原始logit作为$s_i$，不在整个词表上softmax，也不继续自由生成。仅编码器版本对词元表示池化后接线性标量头。前者仍利用预训练解码器的参数和注意计算，后者减少计算并改变聚合方式；不能称二者只是换了最后一个标签词。

对于等级$y_i$，原文比较逐点二分类、成对逻辑损失、列表softmax及其扩展。成对形式为
\begin{equation}
 \mathcal L_{\mathrm{pair}}=
 \sum_{y_i>y_j}\log(1+\exp(s_j-s_i)),
 \label{eq:rs-search-rankt5}
\end{equation}
列表softmax形式为$-\sum_i y_i\log(\exp s_i/\sum_j\exp s_j)$。对非负等级可按协议将$y_i$归一化，使查询之间权重更可控；原文按给定标签加权，这里不把归一化教学变体当成唯一实现。各目标直接作用在候选分数上，梯度更新参与打分的T5参数及标量输出参数。

假设三个候选等级为$(2,1,0)$、分数全为0。逐点二分类若以$y>0$转为相关，只看到$(1,1,0)$，忽略前两者等级差；成对目标包含$A>B,A>C,B>C$三对，总损失$3\log2\approx2.0794$。用归一化等级$(2/3,1/3,0)$作列表目标，模型分布为$(1/3,1/3,1/3)$，交叉熵$\log3\approx1.0986$，分数梯度为$(-1/3,0,1/3)$。两个损失量值不可直接比较优劣，但更新关系可以复算。

线上仍逐对象输出$s_i$后排序，列表关系主要体现在训练目标中。它没有因此选择标题、图片或页面布局，也不会改变硬约束的资格判断。比较monoT5与RankT5，应固定基础检查点、候选、训练数据和输入模板，再分别改变输出接口、损失及结构；跨域收益需在未用于选择模型的集合上检验。

\Needspace{100mm}
% Outline: 9.2.3.11
\subsubsection{RankZephyr：教师监督的列表重排}
\label{sec:rs-search-rankzephyr}

直接调用闭源语言模型得到排列，版本与成本可能限制复现。\term{RankZephyr}把教师的列表次序蒸馏到开放模型：先利用RankGPT3.5对较多训练查询的排列，再用较少RankGPT4排列继续微调。原模型从Zephyr$\beta$的7B检查点出发，教师监督承担相关性比较依据，而不是新增人工评判事实\scite{rs-s23}

输入仍是查询、候选文本及编号，目标是教师生成的编号序列。训练以自回归交叉熵学习目标排列，更新学生参数；筛除格式错误、重复和遗漏的教师输出，加入打乱的输入顺序，并在后续阶段从原20个候选中抽取不同大小的子集，保留教师诱导的相对次序。这样可让模型适应不同窗口长度，也减少只复制初始检索顺序的捷径。

假设教师认为真实对象顺序为$A,C,B$。当输入编号映射为$(1:A,2:B,3:C)$，目标是\texttt{[1]>[3]>[2]}；将输入改为$(1:C,2:A,3:B)$，目标应变为\texttt{[2]>[1]>[3]}。两种输出字符不同，恢复后的对象顺序相同。若学生第二次仍输出\texttt{[1]>[3]>[2]}，得到的是$C,B,A$，说明编号合法并不等于对象排序稳定。训练样本的输入置换必须同步重映射标签。

给定三步教学条件概率0.8、0.7、0.9，教师序列的这三项负对数似然之和为$-\log(0.8\times0.7\times0.9)\approx0.6852$；实际排列还包含括号和分隔等词元，训练时须按真实分词计数。部署继承RankGPT的完整排列核验与滑窗恢复，原文常用窗口20、步长10，也考察多遍重排。多遍结果应连同成倍增加的生成计算报告。

“未用目标集合人工标签”可以描述某种零样本评价条件，但学生确实接受了教师排序监督，不能称为完全没有训练。教师可能延续检索顺序或误读条件，因此验收需要独立相关性标签和输入置换测试，不能只以模仿教师的准确率收尾。固定同一候选池，分别比较教师来源、窗口、输入顺序和推理预算，才能判断改进来自哪里。

\Needspace{100mm}
% Outline: 9.2.3.12
\subsubsection{Rank1：推理计算的相关性重排}
\label{sec:rs-search-rank1}

涉及否定、关系或多项限定时，一次标签输出可能没有足够计算来判断片段是否有用。\term{Rank1}让模型先生成推理文本，再给出\texttt{true}/\texttt{false}判定。其训练从R1教师取得查询—片段推理轨迹，结合MS MARCO正例与困难负例筛选数据，再以序列预测训练学生；原文主要学生来自Qwen2.5基础模型，通过LoRA微调，而非重新联合训练候选检索器\scite{rs-s24}

一条训练样本输入查询与片段，输出由推理词元$\bm r$和末尾判定$y$组成。目标是$-\log p_\theta(\bm r,y\mid q,d_i)$，教师强制下对各输出位置累加交叉熵。筛选会排除教师判定与既有可靠标签不一致的样本，并进一步处理困难负例噪声；这减少部分冲突，但筛选本身也依赖模型判断，不能保证每条推理忠实。

评分需要比“输出了true”更具体。作者实现取得推理末尾两个判定词元的对数概率$\ell_{\mathrm{true}},\ell_{\mathrm{false}}$，再算
\begin{equation}
 s(q,i)=\frac{\exp\ell_{\mathrm{true}}}
 {\exp\ell_{\mathrm{true}}+\exp\ell_{\mathrm{false}}}.
 \label{eq:rs-search-rank1}
\end{equation}
它是给定已生成推理后的二标签分数，不是对所有可能推理轨迹积分的概率。原实现还为未完整结束或未取得两标签概率的输出补做受限判定；部署时应区分自然完成、强制完成与失败回退，不能将三者混为同一种正常预测。

假设$A$说明“密封拉链，适合雨天通勤，当前售价499元”，$B$只说“外观与防水款相似，售价399元”。教学推理分别指出$A$有条件支持、$B$缺防护依据；若末尾二标签概率质量分别为$(0.72,0.08)$与$(0.12,0.48)$，归一化得0.90、0.20。这些分数可以重排已有对象，但价格还须对当前权威字段核验，片段中的陈旧报价不能因模型推理得通而替代实时状态。

增加推理长度会增加每个候选的生成成本。若100个候选平均各生成300词元，总输出就是30000词元，不能只报一次查询的输入长度。应固定候选池，比较有无推理、长度上限、自然停止比例、延迟与独立相关性指标。可读理由也可能是事后解释，甚至在推理中加入原文不存在的事实；理由忠实、条件满足和排序质量需要分别验收，更多计算不自动保证更好。

\Needspace{100mm}
% Outline: 9.2.3.13
\subsubsection{HRNN：查询条件下的历史兴趣排序}
\label{sec:rs-search-hrnn}

用户既浏览过通勤用品也看过登山装备，固定长期画像会把两种兴趣混在一起。\term{HRNN}以会话内循环网络编码查询与先前点击形成短期状态，再以会话间循环网络编码长期历史；查询感知注意从过去会话状态中选择当前有用的部分。本节采用包含该注意的HRNN+QA版本，原工作发表于CIKM 2018，2019年是作者稿上传时间\scite{rs-s38}

设请求前已完成的会话状态为$\bm h_1,\ldots,\bm h_M$，当前查询向量为$\bm q$。由学习函数$a_\psi$得到
\begin{equation}
 \begin{split}
 \alpha_m=\frac{\exp a_\psi(\bm q,\bm h_m)}
 {\sum_{\ell=1}^{M}\exp a_\psi(\bm q,\bm h_\ell)},\qquad\\
 \bm h_{\mathrm{long}}=\sum_m\alpha_m\bm h_m .
 \end{split}
 \label{eq:rs-search-hrnn}
\end{equation}
短期状态$\bm h_{\mathrm{short}}$与长期状态分别投影到文档空间，计算与候选的余弦相似度；再加上公共相关性和历史点击特征经网络形成的分数。原数据不能取得完整公共排序特征，使用原名次作其中一个特征，这个实验条件不应隐去。循环与注意的基本机制复用第~\ref{chap:rs-understanding}~章的用户信息接口及\BookChapterRef{chap:recurrent-networks}{模型卷“循环神经网络”章}。

假设过去两个会话状态简化为通勤$(1,0)$和登山$(0,1)$。含糊查询“背包”的注意权重为$(0.8,0.2)$，长期向量为$(0.8,0.2)$，与通勤、登山候选单位向量的余弦约0.9701、0.2425。查询改为“登山背包”后，权重变为$(0.1,0.9)$，余弦约0.1104、0.9939。若教学公共加短期项分别为0.6、0.5，且长期项权重0.2，两次总分分别为$(0.7940,0.5485)$和$(0.6221,0.6988)$，排序变化可以追溯到查询选择了不同历史。

训练将同一查询中的满意点击与其他候选组成偏好对，以指标加权的成对交叉熵更新循环网络、注意、相似度投影与特征网络，具体目标见第~\ref{sec:rs-search-lambdarank}~节。当前查询还未发生的点击不得进入预测状态；发生后才用于下轮。满意点击也来自当时的曝光集合，Lambda权重强调指标影响，不消除位置与选择偏差。

上线取得已有候选后计算个性化分数，缺历史时退回公共相关性或明示的空历史表示。注意权重高说明模型依赖某段历史，不是用户确认了该意图；跨领域、首次查询及近期意图变化应单独评估。贯穿查询已明确500元上限，历史对高价装备的兴趣不能使600元商品重新取得资格，个性化只能在合格对象间补充比较。

\Needspace{100mm}
% Outline: 9.2.3.14
\subsubsection{KEPS：知识实体增强的个性化排序}
\label{sec:rs-search-keps}

查询中的同名词可能指向不同实体，仅用词向量还不容易说明历史如何消歧。\term{KEPS}根据查询文本与用户长短期历史，为提及的候选实体分配链接概率；通过词与实体的键值记忆形成知识增强画像，再结合当前意图、公共相关性及点击特征给文档评分。它把“此人可能在谈哪个实体”作为排序依据的一部分，并在获得新反馈后更新记忆\scite{rs-s39}

设提及$m$可对应实体集合$E_m$，链接网络产生$a_e(q,H_{u,t})$，概率为
\begin{equation}
 p(e\mid m,q,H_{u,t})=
 \frac{\exp a_e(q,H_{u,t})}{\sum_{e'\in E_m}\exp a_{e'}(q,H_{u,t})}.
 \label{eq:rs-search-keps-linking}
\end{equation}
实体向量按该分布加权形成当前意图，记忆网络根据当前需求读取历史词和实体记录。排名还包含查询、文档的词与实体交互以及知识关系特征，不是只有一项“用户实体向量点乘文档”。知识库实体链接与语言处理衔接第~\ref{part:natural-language-processing}~部分。

按原文的歧义示例，查询“cherry评论”可能涉及花卉、键盘品牌或小说。假设请求前历史使链接分布为$(0.1,0.8,0.1)$，三篇候选分别清楚对应这三种实体，教学公共相关性均为0.5，实体相容贡献权重0.4，则简化总分为$(0.54,0.82,0.54)$。这里的线性式只用于展示链接概率如何改变排序，原网络还使用记忆及多类交互特征。对“防水通勤背包”也可借历史消解含糊品牌名，但不能靠品牌知识猜价格和库存。

训练以会话为单位组织相关与其他候选的成对间隔损失$\max(0,1-s_++s_-)$，共同学习链接、记忆与排名参数。本例若键盘评论是正例、小说评论是负例，损失为$1-0.82+0.54=0.72$。点击发生后，原方法用点击文档中的实体相容信息修正当前链接概率，并在置信条件成立时调整历史提及，再供后续查询使用；更新循环有会话范围限制，不能无限追溯或提前使用未观测反馈。

例如，新点击指向小说，下一轮的教学链接分布更新为$(0.05,0.15,0.80)$，同一简化评分变为$(0.52,0.56,0.82)$。这不是用后来点击重写上一轮预测，而是记录一条新的证据改变后续状态；是否真的应该转向小说，还需排除误点和历史链接错误。实体解释准确率与排序指标应分别评估，AOL日志中的候选和标签条件也不能外推到任意目录。知识覆盖、抽取错误和受偏点击会共同传播，个性化消歧不等于识别了点击的因果效应。

% Outline: 9.2.4
\subsection{结果组织}
\label{sec:rs-search-organization}

相关性次序还不是可用的结果页。同一说明文档的三个相邻片段可能占据前三位，却都回答相同问题；不同商品标题相近，用户也可能看不出为什么价格差了100元。结果组织要让用户识别对象、理解匹配依据、比较差异并完成访问。列表内容与位置选择复用第~\ref{sec:rs-content-organization}~节，表达选择复用第~\ref{sec:rs-presentation-expression}~节，本节补足搜索中的对象和证据要求。

对商品，结果至少区分规格、当前适用价格、库存和防护说明；对文档，保留标题、来源、时间及片段在原文中的位置。标题写“完全防水”而证据仅支持“日常防泼水”，属于展示承诺扩大，即使排序器曾正确读取原文，也会误导用户。生成摘要必须能回到对象的实际资料，不能把查询里要求的属性直接改写成该对象已具备的属性。页面筛选面板也应使用同一字段口径，否则“500元以下”筛选与卡片价格会相互矛盾。

需要覆盖多个方面时，可按第~\ref{sec:rs-mmr}~节等方法在相关性与重复之间权衡。假设前三个高分片段全来自同一份材料介绍，第四个说明通勤携带尺寸，第五个说明防水测试条件；去掉重叠片段、保留不同方面可能更利于探索。但若任务是寻找一份精确手册，机械地增加方面多样性反而可能稀释目标。应先定义任务需要哪些方面，再判断列表是否覆盖，而不是所有查询都强制多样化。

商品、选购文章和测试报告可以分组展示，使“可购买对象”与“辅助决策资料”可辨。一个高相关选购文章不能冒充库存充足的商品；一份检测报告也不应因匹配多个商品而在每个位置重复占满页面。跨来源材料互相矛盾时，应显示来源、适用时间与冲突点，而不是在摘要中悄悄拼成一个没有出处的新结论。

检索增强生成回答位于另一输出接口：检索提供真实片段，生成器组织答案，引用核验检查每项主张是否由所指片段支持。候选命中、回答正确与引用完整性不能互相替代。比如系统取回了正确厂商资料，却在答案中引用另一篇同名企业文章，检索与引用分别得到不同结论；仅生成正确城市名也可能没有足够证据。完整语言生成与RAG流程衔接第~\ref{part:natural-language-processing}~部分，此处保留证据标识、版本和位置供它使用。

广告或促销混入结果页时，沿用第~\ref{sec:rs-ads-allocation}~节与约束接口，同时标明付费身份。支付能力不能代替当前查询相关性，宣传素材不能绕过价格、状态和证据检查；拍卖与收费机制由广告专题展开。最终应检查的不只是前几个对象相关不相关，还包括结果是否可访问、重复是否过多、重要差异是否可辨，以及页面承诺是否与真实对象一致。

% Outline: 9.2.5
\subsection{反馈与检索更新}
\label{sec:rs-search-feedback}

搜索可以在一次返回后完成，也可以因需求或证据变化继续。用户点击、筛选、改写和澄清回答提供关于意图与体验的信息；已经取得的实体、关系和文档则提供关于下一次访问的事实线索。两者应分别记录：用户点了某个结果不代表接受其中全部事实，模型取得一个实体也不代表用户改变了需求。

更新的产物可以是更明确的查询、会话表示、补充候选或停止决定。询问能区分通勤与登山，ConvDR能承接“500元以内的呢”，IRCoT依据前一步材料选择后续查询，HippoRAG通过预建关系取回关联片段。这些方法改变不同接口，既能组合，也能独立使用。原要求清楚且已有充分结果时，额外追问只增加负担。

\Needspace{100mm}
% Outline: 9.2.5.1
\subsubsection{BERT-LeaQuR澄清问题检索}
\label{sec:rs-search-leaqur}

用户只输入“背包”时，系统可以询问用途，但问题本身也要从候选中取得。\term{BERT-LeaQuR}面向既有澄清问题库，分别编码当前查询和每个候选问题，再学习两者是否相关。原文在\term{Qulac}中组织查询主题、需求方面、澄清问题和回答；Qulac是数据集，不能当作这个检索网络的名字\scite{rs-s11}

设候选问题为$a_j$，BERT表示为$\bm h_q,\bm h_{a_j}$，匹配网络$F_\psi$接收两者表示并输出两个logit，经softmax得到$p_j=P(\text{问题相关}\mid q,a_j)$。训练用问题相关标签的交叉熵，微调BERT及前馈匹配网络；原文的隐藏层采用ReLU。它不是对查询与问题直接取一个固定向量内积，虽然两侧表示分别形成。

假设问题库含“主要用于通勤还是登山”“是否需要电脑夹层”“要预订哪个城市的酒店”。对于“背包”，教学模型给分0.85、0.65、0.05，候选预算2时返回前两项。若第一项标签为相关，训练损失为$-\log0.85\approx0.1625$；若第三项标签为不相关，损失为$-\log0.95\approx0.0513$。这两项仅衡量问题是否值得纳入候选，不是回答后文档检索会提高多少。

部署可预计算问题库表示，但匹配网络仍需结合当前查询打分。候选范围影响系统能问什么：问题库没有防护等级问题，模型无法通过排序发明这个候选。库中问题也可能依赖旧上下文或已经被回答，需要在选择阶段结合会话检查。用户明确说“防水通勤背包，价格不超过500元”后，再问“你的预算是多少”通常是重复，问题与主题相关却未必有信息增益。

因此问题候选召回应与最终追问质量分开。取回范围太窄会遗漏有用方面，范围太宽则给下一步增加筛选成本；无合适问题时可以继续展示结果或说明信息不足，不应强行输出排名第一的问题。这里也不能把未来回答作为问题检索特征。

\Needspace{100mm}
% Outline: 9.2.5.2
\subsubsection{NeuQS澄清问题选择}
\label{sec:rs-search-neuqs}

“用于通勤吗”和“是否需要电脑夹层”都可能相关，究竟问哪个，还取决于已经知道什么，以及回答能否改善文档检索。\term{NeuQS}把问题选择作为一个有后续效果的评分任务。其输入包括查询、候选问题、既有问答历史、候选问题所诱导的文档检索分数，以及查询效果预测信号\scite{rs-s11}

原方案用BERT形成查询和问题表示，对已有问答对表示取均值得到历史表示，再与前$k$文档分数及不同截断处的查询效果预测值拼接。两层ReLU前馈网络输出问题分数，采用逐点交叉熵学习。历史为空时须采用固定空状态，而不能计算空集合平均；新候选问题尚未获得回答，线上检索特征只能使用查询、已有历史和该问题本身。

离线Qulac为需求方面提供回答，可以用回答加入检索后的排名效果评估问题质量，并据此组织监督或最优问题对照。这个离线信息不能进入部署特征。假设已知用户回答“用于通勤”，候选问题$A$问电脑尺寸、$B$再次问通勤还是登山；已有历史相同，模型分别输出0.75和0.30，选择$A$。若教学二值质量标签为$(1,0)$，交叉熵之和为$-\log0.75-\log0.70\approx0.6444$。实际质量标签如何由后续指标构造，必须作为训练协议公开。

假设未询问时首个有用文档排第4，倒数排名为0.25；离线该方面对$A$回答“14英寸”后，首个有用文档到第1，指标为1，对$B$重复回答后仍为0.25。这说明为什么问题相关不等于问题有用，也说明标签依赖具体用户方面和检索器。上线选择$A$时不能先知道“14英寸”，只能在用户真正回答后加入条件，重新核验和检索。

评价需按主题与方面划分，避免同一问答模板在训练和测试中重复；问题库覆盖、回答真实性和模拟会话长度都限制结论。最优问题的离线上限只表示存在有帮助的问题，不证明模型总能选到，也不证明每次追问的收益超过交互成本。拒答、含糊回答和误解会累积错误条件，需要允许用户撤销或更正。

\Needspace{100mm}
% Outline: 9.2.5.3
\subsubsection{ConvDR会话稠密检索}
\label{sec:rs-search-convdr}

用户先问“防水通勤背包”，再问“500元以内的呢”，单独编码末轮会丢失对象与用途。\term{ConvDR}将当前查询和会话历史直接编码成检索向量，不要求线上先生成一条完整独立查询。训练时则可以使用人工独立改写的查询作为教师输入，使学生学习在带省略的会话表达下接近教师表示\scite{rs-s12}

设学生$f_\theta$读取当前查询$q_t$及历史$H_t$，教师$g$读取独立改写$q_t^*$，文档侧表示$\bm v_i$来自固定的已有检索器。蒸馏目标为
\begin{equation}
 \mathcal L_{\mathrm{KD}}
 =\frac1d\|f_\theta(q_t,H_t)-g(q_t^*)\|_2^2 .
 \label{eq:rs-search-convdr}
\end{equation}
有少量相关片段时，还可加入基于$f_\theta(q_t,H_t)^\top\bm v_i$的正例对比损失$\mathcal L_{\mathrm{rank}}$。原文分别考察仅蒸馏、仅相关性监督和二者相加的设置；固定教师和文档编码，把梯度用于会话查询编码器，减少重新学习整个空间的需求。

假设教师对“价格不超过500元的防水通勤背包”输出$(0.8,0.6)$，只编码末轮的学生为$(0.1,0.9)$，其二维均方误差为$(0.49+0.09)/2=0.29$；结合历史后学生输出$(0.7,0.6)$，误差为0.005。对通勤片段向量$(1,0)$和泛价格说明向量$(0,1)$，只看末轮得到分数0.1、0.9，结合历史后得到0.7、0.6。这个例子展示历史如何恢复访问方向，价格上限仍由结构化条件另行执行。

使用时编码会话、访问固定文档向量目录并返回真实片段，不必运行教师改写器。教师表示并不提供完美意图：人工改写也可能遗漏条件，旧会话可能含误解或无关话题，长度上限又会截断早期主体。应记录哪些历史被保留、何时重置上下文、用户纠正如何覆盖旧条件。TREC CAsT与OR-QuAC的会话和标签条件不同，不能直接等同真实开放交互中的任务完成。

复杂查询还可能不是省略，而是需要隐含关联。\term{BRIGHT}用论坛中的领域问题及数学、代码任务检验这种检索需求；它是评价基准，不是自动分解查询的算法。改写更充分可以改善表示，但不保证所需的多份证据同时被取回。探索中可以询问、分解需求、补充检索或更换来源，选择依据应是缺少什么；停止则综合用户完成、证据覆盖、边际新增信息和预算，而不是规定必须执行固定轮数\scite{rs-s15}

\Needspace{100mm}
% Outline: 9.2.5.4
\subsubsection{IRCoT：证据条件下的迭代检索}
\label{sec:rs-search-ircot}

有些后续查询只有看过第一份材料才知道该问什么。问题“B17背包制造商的总部在哪”需要先确认制造商，再查其总部。\term{IRCoT}交替执行推理与检索：原问题取得初始片段，语言模型根据已有材料和已生成句子产生下一句推理，再用这句作为查询取得补充片段。原方案利用少样例提示和已有检索器，不额外训练一个统一的端到端模型\scite{rs-s40}

记累计片段集合为$E_\ell$，已生成推理为$r_{1:\ell-1}$。一次循环是
\begin{equation}
 \begin{split}
 r_\ell&=G(q,E_{\ell-1},r_{1:\ell-1}),\\
 E_\ell&=E_{\ell-1}\cup\operatorname{Retrieve}(r_\ell,k).
 \end{split}
 \label{eq:rs-search-ircot}
\end{equation}
即使生成器一次产生多句，原方法每轮只取第一句，以便取得新证据后再继续。若出现规定的答案停止串或达到步数上限，则终止并返回累计片段；还需限制总片段数，使上下文能容纳提示与证据。这个输出集合不同于另一个阅读器生成的最终答案。

表~\ref{tab:rs-search-evidence-steps}给出明确假设的教学证据库，不描述真实厂商。$d_1$记载“B17由澄海制造生产”，$d_2$记载“澄海制造总部位于青州市”，$d_3$仅谈B17外观。
\begin{table*}[htbp]
 \centering
 \begin{tabularx}{\linewidth}{lXXX}
  \toprule
  步骤 & 访问依据 & 本步返回 & 尚缺信息 \\
  \midrule
  初始检索 & 原总部问题 & $d_1,d_3$ & 制造商总部位置 \\
  推理后检索 & “B17制造商是澄海制造” & $d_2,d_1$ & 核对关系与实体同一性 \\
  停止检查 & $d_1$和$d_2$形成关系链 & 去重后$\{d_1,d_2,d_3\}$ & 无新增事实时结束 \\
  \bottomrule
 \end{tabularx}
 \caption{IRCoT教学过程。名称和事实均为假设输入；累计集合包含一份无关片段，不能把返回数当作证据充分性。}
 \label{tab:rs-search-evidence-steps}
\end{table*}

在这个例子中，首轮取2份只覆盖两条必要事实中的1条，第二轮新增$d_2$后覆盖2条。总共两次检索；若第一轮推理产生制造商句，第二轮推理报告答案，则是两次生成调用。与单次检索比较时，既可以固定最终片段预算，也可以固定总访问预算，但必须明确是哪一种；让迭代多花一倍调用后只比较命中，不足以说明相同资源下更好。

中间句若误写为同名的另一家制造商，后续检索可能非常准确地找回错误公司的总部。必须保留句子所依据的片段标识，检查实体与关系，并允许回退或寻找相矛盾来源。停止串只是程序条件，不是证据完备证书；即使最终城市碰巧正确，检索过程仍可能没有取得支持材料。检索覆盖、答案正确、引用支持及推理忠实分别评价，完整RAG与语言行动流程仍衔接第~\ref{part:natural-language-processing}~部分。

\Needspace{100mm}
% Outline: 9.2.5.5
\subsubsection{HippoRAG：关系访问的跨文档证据检索}
\label{sec:rs-search-hipporag}

若大量查询都需要访问跨文档关系，可以先把这些关系组织起来，减少每次逐句生成后续查询的需要。\term{HippoRAG}离线用语言模型从片段抽取实体、名词短语和关系，再用检索编码器连接相近表达；线上将查询实体映射到图节点，通过个性化PageRank传播检索信号，最后汇总回片段。这里“个性化”指查询种子分布，不指用户长期兴趣\scite{rs-s41}

离线先提取实体，再在其辅助下进行开放信息抽取，形成三元组边；编码器相似度超过阈值的近义节点增加连接。还保存节点—片段计数矩阵$\bm B$，$B_{v,i}$为节点$v$在片段$i$中的出现次数。建图使用已有语言模型与检索器，新增索引不等于重新训练它们。抽取结果必须能回溯原文，关系方向、限定时间和否定不能在图中随意丢掉。

查询实体先与节点表示做相似度链接，得到种子集合。原方法还以节点出现于多少片段的倒数作为特异性权重，降低过于常见节点的影响。将加权种子归一化为$\bm b$，设$\bm P$是选定图访问规则下的行随机转移矩阵，传播可写为
\begin{equation}
 \begin{split}
 \bm\pi^{(\ell+1)}&=(1-\alpha)\bm b+
                         \alpha\bm P^\top\bm\pi^{(\ell)},\\
 \bm s&=\bm B^\top\bm\pi,\qquad 0<\alpha<1 .
 \end{split}
 \label{eq:rs-search-hipporag}
\end{equation}
悬空节点需要定义回流规则；迭代到容差或次数上限后，按片段分数$\bm s$取得候选。节点概率归一不意味着片段分数归一，因为一个节点可属于多篇，计数也可能大于1。

复用上一节教学事实，图节点为背包$A$、制造商$F$、城市$C$，两条关系来自$d_1,d_2$。仅为展示传播，令访问图沿$A-F-C$双向走，种子$\bm b=(1,0,0)$、$\alpha=0.8$，从$\bm\pi^{(0)}=\bm b$起步，第一轮为$(0.2,0.8,0)$，第二轮为$(0.52,0.16,0.32)$。若两篇的节点计数分别为$(1,1,0)$与$(0,1,1)$，第二轮片段分数为0.68、0.48，两份证据都获得非零权重。两轮只是可复算中间状态，正式检索仍需继续迭代或明确有限步近似。

图上的双向可达只用于候选访问，不改变事实方向：“厂商生产商品”不能倒推成“商品生产厂商”。即使两篇都被取回，也须核对“B17的制造商”和“该制造商的总部”是否属于同一实体及同一时间。错误近义边可能把两个同名企业连接，常见节点可能扩散到许多无关片段；种子提取遗漏则可能从错误邻域开始。

与IRCoT比较应固定证据库，并计入离线抽取、图存储和在线访问预算。IRCoT以新证据决定下一步查询，HippoRAG以预建图传播信号，两者可以组合，例如在每次IRCoT检索中调用图访问。图路径既不是完整语言推理，也不直接生成可靠答案；最终事实与引用仍需从原始片段验证。

% Outline: 9.3
\section{搜索学习与评价}
\label{sec:rs-search-evaluation}

模型输出看似合理时，仍需知道改善了什么。本节沿需求、候选、次序、页面和交互共同检查，不再新增一个检索环节。训练的相关性等级要绑定查询说明、对象单位和有效时间，困难负例与误负例处理复用第~\ref{sec:rs-retrieval-training}~节；目录与表示更新复用第~\ref{sec:rs-candidate-maintenance}~节。两个近重复片段若同时支持答案，不应只因一个被标注就把另一个当成可靠负例。

查询理解检查原需求和明确条件是否保留，特别关注省略承接、否定和数值；候选检索检查相关且合格对象的覆盖；排序检查已覆盖对象的前置程度；结果页检查方面、来源与可访问性；交互检查任务完成和额外成本。这样才能区分“正确对象没取回”“取回却排低”“排对但显示错误价格”和“追问把原要求改坏”。端到端指标有助于判断用户收益，中间产物有助于定位原因，二者不能只保留其一。

给定已判定的合格相关对象集合$R_q$及去重后的前$k$结果$L_q^{(k)}$，对象召回率为
\begin{equation}
 \operatorname{Recall@}k(q)
 =\frac{|R_q\cap L_q^{(k)}|}{|R_q|}.
 \label{eq:rs-search-recall}
\end{equation}
分母是该查询按固定单位判定的相关对象数，不是目录大小或返回数；$R_q$为空时应明示排除或另外统计无答案查询。若只标注出一部分相关对象，公式计算的是相对当前判断集合的召回，不能直接声称完整相关集合已被覆盖。

已知对象查找更关心首个正确结果。令$r_q^*$为首个满足相关标准的名次，截断外或无命中时贡献0，则
\begin{equation}
 \operatorname{MRR@}k=
 \frac1{|Q|}\sum_{q\in Q}
   \frac{\mathbf1[r_q^*\leq k]}{r_q^*}.
 \label{eq:rs-search-mrr}
\end{equation}
实际实现对无命中直接记0，不以无穷名次参与不明确运算。MRR忽略首个之后的相关对象，所以第一条正确、其他全漏的列表也可以得1，不适合单独评价多份材料覆盖。

分级相关性常采用NDCG。对当前排名第$r$位对象的等级$y_{q,r}$，定义
\begin{equation}
 \begin{split}
 \operatorname{DCG@}k(q)&=
 \sum_{r=1}^{k}\frac{2^{y_{q,r}}-1}{\log_2(r+1)},\\
 \operatorname{NDCG@}k(q)&=
 \frac{\operatorname{DCG@}k(q)}{\operatorname{IDCG@}k(q)}.
 \end{split}
 \label{eq:rs-search-ndcg}
\end{equation}
理想DCG在相同评判对象与等级上按降序计算，零理想增益时需要统一约定。是否把违反硬条件的对象设为0级，或先排除后评价，应事先确定，并另外报告违反率；不能把过滤后的漂亮指标掩盖原始输出中的违规结果。

例如三份合格相关对象等级为3、2、1，另有0级对象。某次前3位等级为$(0,3,2)$，则Recall@3为$2/3$，倒数排名为$1/2$，DCG约为$7/\log_2 3+3/2=5.9165$，理想DCG为$7+3/\log_2 3+1/2\approx9.3928$，NDCG约0.6299。它漏掉一件相关对象、首个命中偏后，且最高等级未在首位；三个指标分别暴露不同问题。

问答文献的“含答案片段命中”常按查询计数：
\begin{equation}
 \operatorname{AnswerHit@}k=
 \frac1{|Q|}\sum_{q\in Q}
 \mathbf1[\text{前}k\text{片段中至少一份含规定答案文字}].
 \label{eq:rs-search-answer-hit}
\end{equation}
分母是查询数，不是全部相关片段数。若要回答制造商总部问题，只含“青州市”三个字的无关片段也可能满足宽松字符串判定，却没有支持关系链。多证据任务还要检查必要事实是否全部由合适来源支持，以及事实之间是否能合法连接；答案文本的准确率和引用主张覆盖率应再单独评价。

未判定结果不能自动解释为已知无关。评测程序可能按惯例把未判定项当0，但报告需注明这一计算口径与标注覆盖，必要时扩充判断池或分别分析已判定子集。同一父文档的重叠片段、镜像页面和近重复商品资料必须按预定规则合并；否则切片更细可能增加命中次数，却没有新增信息。结果页的多个方面也需有方面标注，不能用去重后的对象数代替方面覆盖。

比较协议沿用第~\ref{chap:rs-evaluation}~章的选择质量、时间切分、分群和资源条件：固定目录快照、可用字段、标签、候选截断及预算，隔离查询、来源、时间和近重复。模型参数量相同不意味着资源相同，语言生成次数、上下文长度、近邻访问量、重排候选数和图构建成本都需记录。按查询类型、领域、时间、历史有无和条件复杂度分组，既报告平均值，也检查少数困难组的严重条件违反。

\Needspace{80mm}
\term{BEIR}汇集不同领域与相关性任务，用于检验没有针对目标集合训练时的迁移；各集合相关对象数量、文档长度和标签等级差异很大，同一Recall截断在不同任务中也可能对应不同难度\scite{rs-s13}

前述STaRK增加文本与关系条件，BRIGHT关注需要背景推理才能识别有用材料的查询，FollowIR通过评判指令变化检验模型是否跟随条件。它们覆盖不同失败方式，不能用一个总体均分替代全部检查，也不能把论文时点的平均成绩写成永久算法次序。基准结果应注明版本、标签可见范围及是否仅在预先给定候选中重排，避免把候选池内排序与全目录访问混比。

日志评价需要考虑曝光、选择和反馈解释；随机在线比较与交错实验则必须具有对应的分配机制和归因条件。交错点击胜率不是任意列表差异的无条件因果结论，查询或用户之间的干扰也可能改变估计。用户调查能补充任务完成、理解和满意度证据，但须说明抽样、问题措辞与比较方式，不能因用户主观更喜欢就宣称随机干预收益。最终结论应同时保留任务效果、交互负担和适用范围。

% Outline: 9.3.1
\subsection{CIRR：组合图像检索评价基准}
\label{sec:rs-search-cirr}

组合图像检索的标签既要求改变指定部分，又隐含保留未改动的部分。\term{CIRR}以真实图片、参考—目标图对和人工相对描述形成实例，并从视觉相近图片中构造小型子集，检验细粒度选择。CIRPLANT是同文提出的模型，CIRR则固定输入、候选范围与评价条件，不能把数据集成绩误称为另一模型的输出\scite{rs-s37}

原构造先形成6张相似图像的子集，再从中选择参考—目标对，让标注者在看到其他图像的条件下编写能够区分目标的修改句。评测排除参考图，因此子集检索通常在其余5张中比较；全目录评价则在对应划分的图像目录中找目标。两者都可报告Recall@$k$，但一个查询只有指定目标时，它实际上是目标是否进入前$k$的命中指标，再跨查询平均，不是对业务中所有合法目标的完整召回率。

假设参考图是某轮廓的红色背包，文字为“保持轮廓，改成蓝色”。指定目标$A$是同轮廓蓝包，$B$是另一张也满足要求但未被指定的图片，$C$是背景相似却轮廓不同的蓝包。设全目录次序为$C,B,D,A,\ldots$，$A$名次4，因此该查询Recall@1为0、Recall@5为1。若相似子集不含$B,D$而包含$C,A$及另外三图，目标在子集中排2，子集Recall@1为0、Recall@2为1。候选范围一变，数字就不能直接横比。

这个例子中$B$可能是业务合法结果，基准却只按指定$A$判断。不能由此宣布基准无用，也不能把$B$无条件当作错误：应保留基准原协议，同时用补充标注诊断多合法目标。原数据另有辅助标注帮助解释主体、背景和隐含条件，原论文并未在其模型训练中使用全部这些标注；引用时不能把可用资料写成模型实际监督。

固定图像划分、相对描述、参考图排除规则、子集与全目录范围以及预算，才能比较TIRG、CIRPLANT或后续方法。图像相似不替代商品资格：即使检索到了正确颜色与轮廓，还需通过第~\ref{sec:rs-search-candidate-set}~节核实价格、库存和规格。背景捷径、主体指向错误与未标合法目标应分开统计，较小子集的高分也不能等同大目录里的访问效果。

% Outline: 9.4
\section{本章小结}
\label{sec:rs-search-summary}

显式查询给共同架构增加了一个必须保持一致的依据：当前需求。查询理解将它变成可执行表达与条件，候选接口从当前目录取得真实对象，相关性排序比较对象满足需求的程度，结果组织帮助用户识别、比较和访问。个性化可以消歧和补充情境，却不能覆盖明确价格上限，也不能把历史兴趣当成这次需求的最终裁决。

词项、稠密、多向量和标识生成的差异，分别落在表示、匹配、监督与访问接口上。选择方法需要同时考虑精确条件、词汇差异、训练资料、目录更新和计算预算，不能只按模型新旧排列优劣。硬条件核验、对象映射与来源记录始终保留，生成的查询、标识、相关标签和答案则对应不同产物，不能相互冒充。

需求不清时可以询问，表达省略时可以承接历史，证据缺失时可以迭代检索或访问关系图；已有充分结果时也可以停止。评价要分别检查覆盖、次序、条件违反、页面承诺与任务完成，并区分答案文字命中、证据充分和最终回答正确。这样才能知道系统究竟在哪一步帮助了用户，又在哪些信息、目录或反馈条件下仍会失败。

当结果页包含有偿曝光时，系统还需处理出价、支付和预算，但付费对象仍应满足用户的显式需求。下一章由此转向广告中的选择、收费与效果识别。
